\documentclass[12pt]{amsart}
\usepackage{amsthm,amsfonts,amssymb,amscd,mathrsfs}
\usepackage{bezier,longtable,amstext}
\usepackage{amssymb}

\sloppy \topmargin-10mm \oddsidemargin-5mm \evensidemargin-5mm

\setlength{\textwidth}{17cm} \setlength{\textheight}{25cm}
\usepackage[all]{xy}
%\usepackage[pagebackref,colorlinks]{hyperref}
%\usepackage{showkeys}


\newcommand{\EXT}{{\mathcal Ext}}
\newcommand\uC{{\underline{C}}}
\newcommand\Aut{\text{Aut}}
\newcommand\codim{\text{codim}}
\newcommand\coker{\operatorname{coker}\nolimits}
\newcommand{\const}{\operatorname{const}\nolimits}
\newcommand\ev{\text{ev}}
\newcommand\ext{\text{ext}}
\newcommand\Ext{\operatorname{Ext}\nolimits}
\newcommand\GL{\text{GL}}
\newcommand{\Hilb}{\operatorname{Hilb}\nolimits}
\newcommand\hhom{\text{hom}}
\newcommand\Hom{\operatorname{Hom}\nolimits}
\newcommand\id{\text{id}}
\newcommand\im{\operatorname{Im}\nolimits}
\newcommand\Ker{\operatorname{Ker}\nolimits}
\newcommand\length{\text{length}}
\newcommand\mult{\text{mult}}
\newcommand\odd{\text{odd}}
\newcommand{\Pic}{\operatorname{Pic}\nolimits}
\newcommand\pr{\text{pr}}
\newcommand{\Proj}{\operatorname{Proj}\nolimits}
\newcommand\red{\text{red}}
\newcommand\res{\text{res}}
\newcommand\rk{\text{rk}}
\newcommand\Fl{\text{Fl}}
\newcommand{\Sing}{\operatorname{Sing}\nolimits}
\newcommand\Span{\text{Span}}
\newcommand\Stab{\text{Stab}}
\newcommand\Sym{\operatorname{Sym}\nolimits}
\renewcommand{\phi}{\varphi}

\newcommand{\refle}[1]{Lemma \ref{#1}}
\newcommand{\refth}[1]{Theorem \ref{#1}}
\newcommand{\refcor}[1]{Corollary \ref{#1}}
\newcommand{\refprop}[1]{Proposition \ref{#1}}
\newcommand{\refsec}[1]{Section~ \ref{#1}}
\newcommand{\refeq}[1]{(\ref{#1})}

\newcommand {\eqdef} {:=}
\newcommand {\eqdefrev} {=:}

\newcommand\elm{{\operatorname{elm}\nolimits}}
\newcommand\Supp{{\operatorname{Supp}\nolimits}}

\newtheorem{theorem}{Theorem}[section]
\newtheorem{proposition}[theorem]{Proposition}
\newtheorem{prop}[theorem]{Proposition}
\newtheorem{conjecture}[theorem]{Conjecture}
\newtheorem{lemma}[theorem]{Lemma}
\newtheorem{corollary}[theorem]{Corollary}
\newtheorem{sub}[subsection]{}
\newtheorem{subsub}{}[subsection]
\newtheorem{subsubsub}{}[subsubsection]
\newtheorem{assumption}[theorem]{Assumption}
\newtheorem{claim}[theorem]{Claim}
\newtheorem{notation}[theorem]{Notation}



\theoremstyle{definition}
\newtheorem{definition}[theorem]{Definition}
\newtheorem{example}[theorem]{Example}
\newtheorem{proposition-definition}[theorem]{Proposition-Definition}

\newtheorem{remark}[theorem]{Remark}

\begin{document}


\title{Triviality of vector bundles on twisted ind-Grassmannians}

\author[I.Penkov]{\;Ivan~Penkov}

\address{
Jacobs University Bremen \\
School of Engineering and Science, Campus Ring 1, 28759 Bremen, Germany}
\email{i.penkov@jacobs-university.de}

\author[Tikhomirov]{\;Alexander S.~Tikhomirov}


\address{
Department of Mathematics\\
State Pedagogical University\\
Respublikanskaya Str. 108
\newline 150 000 Yaroslavl, Russia}
\email{astikhomirov@mail.ru}







\maketitle


\thispagestyle{empty}

\begin{abstract}
Twisted ind-Grassmannians are ind-varieties $\mathbf{G}$ obtained as direct limits of
Grassmannians $G(i_m,V^{n_m})$ for $m\in\mathbb{Z}_{>0}$, under embeddings
$\phi_m:G(i_m,V^{n_m})\to G(i_{m+1}, V^{n_{m+1}})$ of degree greater than one. It has
been conjectured by J.~Donin and I.~Penkov \cite{DP}, and by I.~Penkov and  A.~Tikhomirov
\cite{PT}, that any vector bundle of finite rank on a twisted ind-Grassmannian is trivial.
We prove this conjecture.

2000 Mathematics Subject Classification, Primary 14M15 (Secondary 14J60, 32L05).

\keywords{Keywords: ind-varieties, twisted ind-Grassmannians, vector bundles.}

%\thanks{Partially supported by grant DFG PE980/3-1.}
\end{abstract}


\markright{Triviality of vector bundles on twisted ind-Grassmannians}


%\thanks



\section{Introduction and statement of the main result}\label{sec1}
\label{Introduction} An {\it ind-Grassmannian} $\mathbf{G}=\underset{\to}\lim
G(i_m,V^{n_m})$ is an ind-variety obtained as the direct limit of a chain of
embeddings
\begin{equation}\label{eq1}
G(i_1,V^{n_1})\stackrel{\phi_1}{\hookrightarrow}G(i_2,V^{n_2})
\stackrel{\phi_2}{\hookrightarrow}\dots\stackrel{\phi_{m-1}}{\hookrightarrow}G(i_m,V^{n_m})
\stackrel{\phi_m}{\hookrightarrow}\dots,
\end{equation}
where $G(i,V)$ denotes the Grassmanian of $i$-dimensional subspaces in a finite
dimensional vector space $V$. Each embedding $\phi_m$ has a well defined degree $\deg
\phi_m$, and the ind-Grassmannian $\mathbf{G}$ is {\it twisted} if $\deg\phi_m>1$ for
infinitely many $m$. A {\it vector bundle $\mathbf{E}$ of rank}
$\mathbf{r}\in\mathbb{Z}_{>0}$ on $\mathbf{G}$ is the inverse limit
$\underset{\leftarrow}\lim E_m$ of an inverse system of vector bundles $E_m$ or rank
$\mathbf{r}$ on $G(i_m,V^{n_m})$ (i.e. a system of vector bundles $E_m$ with fixed
isomorphisms $\psi_m: E_m\cong \phi^*_mE_{m+1}$).

In the special case when $i_m=1$ and $\deg\phi_m=1$ for all $m$,  the study of finite
rank vector bundles on ind-Grassmannians goes back to W. Barth, A. Van de Ven and A.
N. Tyurin, \cite{BV}, \cite{T}. In this case $\mathbf{G}$ is just the infinite
projective space $\mathbf{P}^\infty$, and the remarkable Barth-Van de Ven-Tyurin
Theorem claims that any vector bundle of finite rank on $\mathbf{P}^\infty$ is
isomorphic to a direct sum of line bundles. Historically this is the first
manifestation of a general phenomenon that seems to take place for ind-varieties
defined via sequences of embeddings similar to (\ref{eq1}) with $G(i_m,V_m)$ replaced
by arbitrary flag varieties: in all cases known, the restriction of any finite rank
vector bundle on the ind-variety to a large enough finite-dimensional flag subvariety
is equivariant. Around the same time this phenomenon occured also in the work of E.
Sato who gave an independent proof of the Barth-Van de Ven-Tyurin Theorem, \cite{S1}.
Shortly after that Sato established a more general result which applies in particular
to the ind-Grassmannian $\mathbf{G}(i,V)$ of $i$-dimensional subspaces in a
countable-dimensional vector space $V$ \cite{S2}.

More recently the subject has been revisited in the papers \cite{DP}, \cite{CT} and
\cite{PT}. In particular, in \cite{PT} a general conjecture about finite rank vector
bundles on ind-Grassmannians $\mathbf{G}$ has been stated. In fact, as we show in
\cite{PT}, if $\mathbf{G}$ is not twisted (which is equivalent to assuming that
$\deg\phi_m=1$ for all $m$), this conjecture is a relatively straightforward corollary
of Sato's result. This leaves open the case of a twisted ind-Grassmannian
$\mathbf{G}$, where the conjecture claims simply that a finite-rank vector bundle on
$\mathbf{G}$ is trivial. So far this latter conjecture was established in the
following three cases: for a rank-two bundle on any twisted ind-Grassmannian
\cite{PT}, for any finite-rank bundle on any twisted projective ind-space (a twisted
projective ind-space can be defined via the sequence (\ref{eq1}) for $i_m=1$ and
$\deg\phi_m>1$ for all $m$) \cite{DP}, and for an arbitrary finite-rank bundle on some
special twisted ind-Grassmannians (for which the embeddings $\phi_m$ are twisted
extensions as defined in \cite{DP}).

In the present paper we prove the conjecture, i.e. the following theorem.

\begin{theorem}\label{main}
A finite-rank vector bundle $\mathbf{E}=\underset{\leftarrow}\lim~E_m$ on any twisted
ind-Grassmannian $\mathbf{G}=\underset{\rightarrow}\lim~G(i_m,V^{n_m})$ is trivial.
\end{theorem}

Here is a brief description of the main ingredients in the proof of Theorem
\ref{main}. First of all, without loss of generality we can assume that $\mathbf{E}$
is self-dual. This is achieved by possibly replacing of $\mathbf{E}$ with
$\mathcal{E}nd\mathbf{E}$. The ultimate goal of the proof is to construct, for large
$m$, subsheafs $\mathcal{F}_m$ of the vector bundles $E_m$ with $c_1(\mathcal{F}_m)>0$
under the assumption that $E_m$ is nontrivial. This then easily leads to a
contradiction since the facts that $\mathbf{G}$ is twisted and $E_m$ is infinitely
extendable force $c_1(\mathcal{F}_m)$ to be infinite. The general idea of such a
construction goes back to Barth-Van de Ven and Tyurin in the case of
$\mathbf{P}^\infty$.

The construction of $\mathcal{F}_m$ combines several ideas and is based on a study of
the variety of maximal jumping lines of the vector bundle $\mathbf{E}$. In our case we
investigate the variety of maximal jumping lines of $E_m$ on $G(i_m,V^{n_m})$. We
reduce the problem to the study of a similar variety for projective space by using a
birational isomorphism of $G(i_m,V^{n_m})$ with a fibred space $X_m$ with fibre a
projective space. A key result in this connection is the existence of universal bounds
for the degree and codimension of the variety of maximal jumping lines through a point
of a vector bundle on a projective space.

The paper is organized as follows. Section \ref{bound in Pn} is a study of varieties
of bounded degree and codimension in projective spaces of growing dimension. The main
result here is that any two points of such a variety can be connected by chain of
projective subspaces of growing dimension. This result is close in spirit to a
classical result of A. Predonzan, and is part of the present paper due to the lack of
a suitable reference.

In section \ref{m-reg} we give a sufficient condition on an integer $m$ for a given
vector bundle $E$ on $\mathbb{P}^n$ to be $m$-regular in the sense of
Mumford-Castelnuovo, i.e. that $H^i(E(m-i))=0$ for $i\ge1$. This condition on $m$ is
needed for the estimate of the degree of the variety of maximal jumping lines through
a point of a vector bundle on a projective space, given in section \ref{bounded
degree}. This estimate (see Theorem \ref{restricted degree}) is given in terms of
rank, second Chern class, maximal jump and dimension of the projective space, under
the assumption that the first Chern class vanishes.

The final section \ref{proof} is devoted to the construction of the subsheaf
$\mathcal{F}_m$ of $E_m$, where $\mathbf{E}=\underset{\leftarrow}\lim~E_m$ is a
self-dual vector bundle on $\mathbf{G}$. Here we replace $G(i_m,V^{n_m})$ by a fibred
space $X_m$, to the fibres of which we apply all above results on vector bundles on
projective spaces. The construction of $\mathcal{F}_m$ then quickly leads to a
contradiction with the nontriviality of $E_m$ as explained above.

We conclude this introduction with an example of a twisted ind-Grassmannian for
which our theorem provides a nontrivial statement.
In this example a twisted ind-Grassmannian arises naturally as a homogeneous space
of a locally linear ind-group. Various further examples of twisted
ind-Grassmannians can be found in the papers \cite{DP} and \cite{PT}.

An interesing ind-group is the ind-group $\mathbf{SL}(n,{\rm Adj})$. Fix $n$ and
consider the embedding
$$
\ SL(n)\to SL(n^2-1)
$$
defined by the requirement that the natural representation of $SL(n^2-1)$ becomes the
adjoint representation when restricted to $SL(n)$. Setting $\mathcal{G}_1:=SL(n),\
\mathcal{G}_2:=SL(n^2-1)$, and iterating this construction we obtain the ind-group
$\mathbf{SL}(n,{\rm Adj})$ as the direct limit $\underset{\to}\lim\ \mathcal{G}_m$.
Fix a subspace $V_1\subset\mathbb{C}^n$. Then
$V_2:=V_1\otimes(\mathbb{C}^n/V_1)^{\vee}$ ( ${}^\vee$ indicates dual space) is a
well-defined subspace of the adjoint representation of $\mathcal{G}_2$.
Iteration of this construction yields a subspace $V_m$ of the natural
representation of $\mathcal{G}_m$ for each $m$. The stabilizers $\mathcal{P}_m\subset
\mathcal{G}_m$ of the spaces $V_m$ form a direct system of parabolic subgroups
$\mathcal{P}_1\subset \mathcal{P}_2\subset...$ with the property $\mathcal{P}_m\cap
\mathcal{G}_{m-1}=\mathcal{P}_{m-1}$. This defines closed embeddings
$\zeta_{m-1}:\mathcal{G}_{m-1}/\mathcal{P}_{m-1}\hookrightarrow
\mathcal{G}_m/\mathcal{P}_m$, and hence an ind-variety
$\underset{\to}\lim\ \mathcal{G}_m/\mathcal{P}_m$. Since
each $\mathcal{G}_m/\mathcal{P}_m$ is a Grassmannian,
$\underset{\to}\lim\ \mathcal{G}_m/\mathcal{P}_m$
is an ind-Grassmannian. Moreover, the restriction of
the tautological bundle on $\mathcal{G}_m/\mathcal{P}_m$ to
$\mathcal{G}_{m-1}/\mathcal{P}_{m-1}$ is isomorphic the cotangent bundle of
$\mathcal{G}_{m-1}/\mathcal{P}_{m-1}$. This shows that the degree of $\zeta_m$ equals
the dimension of the natural representation of $\mathcal{G}_m$. Hence
$\underset{\to}\lim\ \mathcal{G}_m/\mathcal{P}_m$ is a twisted ind-Grassmannian. It
is an exercise to check that the ind-group
$\underset{\to}\lim \mathcal{P}_m$ has no non-trivial finite-dimensional
representations. Therefore $\underset{\to}\lim\ \mathcal{G}_m/\mathcal{P}_m
=\mathbf{SL}(n,{\rm Adj})/(\underset{\to}\lim \mathcal{P}_m)$ admits
no non-trivial $\mathbf{SL}(n,{\rm Adj})$-equivariant vector bundles of finite rank.
Theorem \ref{main}, however, yields the much stronger result that any finite rank
vector bundle on $\underset{\to}\lim\ \mathcal{G}_m/\mathcal{P}_m$ is trivial.


\textbf{Acknowledgements.} We thank F. L. Zak for pointing out to us the work of R.
Braun and S. M\"uller-Stach. We acknowledge the support and hospitality of the Max
Planck Institute for Mathematics in Bonn where the present paper was conceived. This paper
was completed at Jacobs University Bremen and A. S. T. has been partially supported by ICTS at Jacobs University Bremen, as well as by DFG grant PE980/3-1.


\section{Notation and Conventions}\label{sec2}
Our notation is mostly standard. The ground field is $\mathbb{C}$. All vector bundles
considered are assumed to have finite rank. We do not make a distinction between
locally free sheaves of finite rank and vector bundles. We use the term {\it algebraic
variety} or simply {\it variety} as shorthand for a reduced noetherian scheme.
If $\mathcal{F}$ is a sheaf of
$\mathcal{O}_X$-modules on an algebraic variety or a scheme $X$, $\mathcal{F}^j$
denotes the direct sum of $j$ copies of $\mathcal{F}$, $H^j(\mathcal{F})$ denotes the
$j^{th}$ cohomology group of $\mathcal{F}$, $h^j(\mathcal{F}):=\dim H^j(\mathcal{F})$,
and $\mathcal{F}^\vee$ stands for the dual sheaf, i. e. $\mathcal{F}^\vee:=
\mathcal{H}om_{\mathcal{O}_X}(\mathcal{F},\mathcal{O}_X)$. $Sym^j$ and $\wedge^j$
denote respectively j-th symmetric and exterior power. If $Z\subset X$ is a
subvariety, $\mathcal{I}_{Z,X}$ denotes the sheaf of ideals corresponding to $Z$. By
$\mathbb{P}(E)$ we denote the projectivization of a vector bundle $E$ (in particular,
of a vector space).

By a \emph{projective subspace} $\mathbb{P}^k$ in $G(i,V)$ we mean linearly embedded
projective subspace, i.e. the set of $i$-dimensional subspaces $W$ of $V$ with
$V_0\subset W\subset V_1$, where $V_0\subset V_1$ is a fixed flag of subspaces of $V$
of dimensions $i-1$ and $i+k$, or $i-k$ and $i+1$ respectively. In particular, a
\emph{line} in $G(i,V)$ is determined by a flag $V_1\subset V_2$ of subspaces in $V$
with $\dim V_1=i-1$, $\dim V_2=i+1$.

If $C\subset X$ is a smooth irreducible rational curve in an algebraic variety $X$ and
$E$ is a vector bundle on $X$, then by a classical theorem of Grothendieck,
$\displaystyle E|_C$ is isomorphic to $\bigoplus_i\mathcal{O}_C(\delta_i)$ for some
$\delta_1\geq \delta_2\geq\dots\geq \delta_{\rk E}$. We call the ordered $\rk E$-tuple
$(\delta_1,\dots,\delta_{\rk E})$ \emph{the splitting type} of $E|_C$.

Let $E$ be a vector bundle on $G(i,V)$. For an arbitrary rational curve $C$ in
$G(i,V)$ consider the splitting type $(\delta_1,\dots,\delta_{\rk E})$ of the bundle
$E|_C$ and set
$$
\delta_A(E|_C):=\delta_1,\ \ \ \delta_B(E|_C):=\delta_{\rk E},\ \ \
\delta(E|_C):=\delta_A(E|_C)-\delta_B(E|_C),
$$
$$
\kappa_A(E|_C):=\max\{k|1\le k\le\rk E,\ \delta_k=\delta_A(E|_C)\}.
$$
Furthermore, set
$$
\delta_A(E):=\underset{l}\max~ \delta_A(E|_l),\ \ \ \ \delta_B(E):=\underset{l}\min~
\delta_B(E|_l),\ \ \ \
$$
where $l$ runs over all lines in $G(i,V)$,
$$
\delta(E):=\delta_A(E)-\delta_B(E),
$$
$$
\kappa_A(E):= \max\{ \kappa_A(E|_l)\ |\ l\ {\rm is\ a\ line\ in}\ G(i,V)\ {\rm such\
that}\ \delta_A(E|_l)=\delta_A(E)\}.
$$
It is essential to note that $\delta_A(E|_C)$ and $\kappa_A(E|_C)$ are semicontinuous
functions of $C$, where $C$ belongs to any fixed flat family of rational curves in
$G(i,V)$ \cite[Ch. III, Thm. 12.8]{H}.

We need also a notation concerning polynomials. For an arbitrary nonzero polynomial
$f(y_1,...,y_q)=\sum \frac{a_{i_1...i_q}}{b_{i_1...i_q}}y_1^{i_1}...y_q^{i_q}\in
\mathbb{Q}[y_1,...,y_q]$ with coprime $a_{i_1...i_q}\in\mathbb{Z}$ and
$b_{i_1...i_q}\in\mathbb{Z}$ for all $i_1,...,i_q$, we denote by
$f(y_1,...,y_q)^+\in\mathbb{Z}[y_1,...,y_q]$ the polynomial $\sum
a^2_{i_1...i_q}y_1^{2i_1}...y_q^{2i_q}$. Note that $-f(y_1,...,y_q)^+\le
f(y_1,...,y_q)\le f(y_1,...,y_q)^+$ for all $y_1,...,y_q\in\mathbb{Z}$.


\vspace{1cm}

\section{Projective subspaces in varieties of bounded codimension and degree}\label{bound in Pn}
\vspace{0.5cm}

In this section we prove that any two points of a subvariety of bounded codimension
and degree in a projective space of growing dimension can be connected by a chain of
projective subspaces of growing dimension lying on this subvariety. This is a chapter
of the theory of Fano schemes in the spirit of Altman and Kleiman \cite{AK}, and is
also close to Predonzan's Theorem (1948), a modern presentation of which can be found
in \cite{BM}. Thoughout the section $d\in\mathbb{Z}_{\ge2}$ is fixed and
$n\in\mathbb{Z}_{\ge6}$ is variable. The integer $k\in\mathbb{Z}_{\ge1}$ is variable
and satisfies
\begin{equation}\label{k(n-k)>...}
n\ge d\binom{k+d}{d}+k,
\end{equation}
for instance, one may set $k=k(n):=\big[\sqrt[d+1]{n/d}\ \big]$.
\begin{sub}{\bf Projective subspaces in hypersurfaces of bounded degree and growing dimension.}
\label{hypersurf} \rm Consider the projective space $\mathbb{P}^n=\mathbb{P}(V)$ where
$V$ is a vector space of dimension $n+1$. Let
$$
\mathbb{P}^s:=|\mathcal{O}_{\mathbb{P}^n}(d)|,\ \ \ s=\binom{n+d}{d}-1,
$$
be the complete linear series of hypersurfaces of given degree $d$ in $\mathbb{P}^n$.
Consider the natural diagram
\begin{equation}\label{p,q}
G(k+1,V)\overset{\tilde{p}}\leftarrow\Gamma\overset{\tilde{q}}\rightarrow\mathbb{P}^s,
\end{equation}
where $\Gamma=\{(\mathbb{P}^k,H)\in G(k+1,V)\times\mathbb{P}^s\ |\ \mathbb{P}^k\subset
H\}$ and we interpret $G(k+1,V)$ as the Grassmannian of k-dimensional projective
subspaces in $\mathbb{P}^n$. For each pair $(\mathbb{P}^k,H)\in\Gamma$ choose
homogeneous coordinates $(x_0:x_1:...:x_n)$ in $\mathbb{P}^n$ such that
$\mathbb{P}^k=\{x_{k+1}=...=x_n=0\}$. Let $H=\{f(x_0,...,x_n)=0\},\ \ f\in
H^0(\mathcal{O}_{\mathbb{P}^n}(d))$, and
$$
\Phi_i(x_0,x_1,...,x_k):=\frac{\partial f}{\partial
x_{k+i}}(x_0,x_1,...,x_k,0,...,0),\ \ \ 1\le i\le n-k.
$$
Assume that $H$ is smooth. Then
$\overset{n-k}{\underset{i=1}\cap}\{\Phi_i(x_0,x_1,...,x_k)=0\}=\emptyset$ and we have
an exact sequence of normal bundles on $\mathbb{P}^k$
\begin{equation}\label{normal bdls}
0\to N_{\mathbb{P}^k/H}\to \mathcal{O}_{\mathbb{P}^k}(1)^{n-k}
\overset{\epsilon_k}\to\mathcal{O}_{\mathbb{P}^k}(d)\to0,\ \ \ \ \ \
\epsilon_k=(L_1,...,L_{n-k})\mapsto\sum_{i=1}^{n-k}L_i\Phi_i.
\end{equation}
Assume $H$ is generic in the sense that
\begin{equation}\label{SpanPhi=}
\Span(\Phi_1,...,\Phi_{n-k})=H^0(\mathcal{O}_{\mathbb{P}^k}(d-1)).
\end{equation}
Then the exact sequence obtained from (\ref{normal bdls}) via twisting by
$\mathcal{O}_{\mathbb{P}^k}(-1)$ induces a surjective homomorphism
$H^0(\mathcal{O}_{\mathbb{P}^k}^{n-k})\to H^0(\mathcal{O}_{\mathbb{P}^k}(d-1))$, and
it is easy to see that, after twisting back by $\mathcal{O}_{\mathbb{P}^k}(1)$, we get
a surjective homomorphism $h^0(\epsilon_k):H^0(\mathcal{O}_{\mathbb{P}^k}(1)^{n-k})\to
H^0(\mathcal{O}_{\mathbb{P}^k}(d))$. Therefore

\begin{equation}\label{h0(N)}
h^0(N_{\mathbb{P}^k/H}) =(k+1)(n-k)-\binom{k+d}{d}>0,\ \ \
h^1(N_{\mathbb{P}^k/H})=h^1(N_{\mathbb{P}^k/H}(-1))=0
\end{equation}
(the inequality follows from (\ref{k(n-k)>...})).

Note that $\tilde{p}:\Gamma\to G(k+1,V)$ is a projective bundle with fibre
$\mathbb{P}(H^0(\mathcal{I}_{\mathbb{P}^k,\mathbb{P}^n}(d)))$, hence $\Gamma$ smooth
and irreducible. Therefore $\dim \tilde{q}^{-1}(H)\ge\dim\Gamma-s=\dim G(k+1,V)+
\dim\mathbb{P}(H^0(\mathcal{I}_{\mathbb{P}^k,\mathbb{P}^n}(d)))-s=
(k+1)(n-k)+(s-\binom{k+d}{d})-s=(k+1)(n-k)-\binom{k+d}{d}$. From this and
(\ref{h0(N)}) we obtain by deformation theory that
$$
B_H:=\tilde{p}(\tilde{q}^{-1}(H))
$$
has dimension
\begin{equation}\label{dim BX}
\dim B_H=h^0(N_{\mathbb{P}^k/H})=(k+1)(n-k)-\binom{k+d}{d}
\end{equation}
and is smooth at the point $\mathbb{P}^k$ for a generic smooth $H\in\mathbb{P}^s$.
Moreover, the projective morphism $\tilde{q}$ is dominant. Since the image of a
projective morphism is closed \cite[Ch. II, \S4, Thm. 4.9]{H}, this implies that
$\tilde{q}$ is surjective.

\begin{lemma}\label{BXirreducible}
For a smooth generic (in the sense of (\ref{SpanPhi=})) hypersurface
$H\in\mathbb{P}^s$, $B_H$ is a smooth irreducible variety of dimension
$(k+1)(n-k)-\binom{k+d}{d}$.
\end{lemma}
\begin{proof}
The smoothness of $B_H$ follows from the fact that $B_H$ is a generic fibre of the
surjective morphism $\tilde{q}:\Gamma\to\mathbb{P}^s$ of smooth varieties \cite[Ch.
III, \S10, Cor.10.7]{H}.

Let $S_{k+1}$ be the rank-$(k+1)$ tautological bundle on $G(k+1,V)$. By \cite[Thm.
1.3]{AK} $B_H$ is the zero-scheme of a regular section $\sigma\in H^0(T^\vee)$, where
$T:=Sym^dS_{k+1}$. Moreover, we have the standard Koszul resolution of the sheaf
$\mathcal{O}_{B_H}$
\begin{equation}\label{resolutionOBX}
0\to\wedge^{\rk T}T\to...\to\wedge^2T\to
T\overset{\sigma^\vee}\to\mathcal{O}_{G(k+1,V)}\to \mathcal{O}_{B_H}\to0.
\end{equation}
We will show that
\begin{equation}\label{vanishHj}
H^0(T)=H^j(\wedge^jT)=0,\ \ \ 1\le j\le \rk T.
\end{equation}
For this, consider the incidence diagram
\begin{equation}\label{G(i+1,V)}
G(i+1,V)\overset{p_i}\leftarrow Fl(i,i+1,V)\overset{q_i}\rightarrow G(i,V),\ \ \ 1\le
i\le k.
\end{equation}
On $Fl(i,i+1,V)$ one has an exact sequence of vector bundles
\begin{equation}\label{taut}
0\to q_i^*S_i\overset{\theta_i}\to p_i^*S_{i+1}\to q_i^*\mathcal{O}_{G(i,V)}(1)\otimes
p_i^*\mathcal{O}_{G(i+1,V)}(-1)\to0,
\end{equation}
where $S_i$ be the rank-$i$ tautological bundle on $G(i,V)$. Restricting (\ref{taut})
to a fibre $\mathbb{P}^{n-i}_y:=q_i^{-1}(y)$ for $y\in G(i,V)$, we obtain an exact
triple
$$
0\to  q_i^*S_i|_{\mathbb{P}^{n-i}_y}\overset{\theta_i|_{\mathbb{P}^{n-i}_y}}\to
p_i^*S_{i+1}|_{\mathbb{P}^{n-i}_y}\to \mathcal{O}_{\mathbb{P}^{n-i}_y}(-1)\to0,\ \ \ \
\ \ q_i^*S_i|_{\mathbb{P}^{n-i}_y}\simeq (\mathcal{O}_{\mathbb{P}^{n-i}_y})^i.
$$
Passing to symmetric powers and setting $s_i:=\binom{d+i}{d-1},\
t_i:=\binom{d+i-1}{d}$, we have
\begin{equation}\label{sym powers1}
0\to  q_i^*Sym^dS_i|_{\mathbb{P}^{n-i}_y}\to
p_i^*Sym^dS_{i+1}|_{\mathbb{P}^{n-i}_y}\to
\overset{s_i}{\underset{p=1}\oplus}\mathcal{O}_{\mathbb{P}^{n-i}_y}(a_p)\to0,\ \ \
-d\le a_p\le-1,\ \ 1\le p\le s_i,
\end{equation}
\begin{equation}\label{sym powers2}
q_i^*Sym^dS_i|_{\mathbb{P}^{n-i}_y}\simeq (\mathcal{O}_{\mathbb{P}^{n-i}_y})^{t_i} \ .
\end{equation}

Consider the exact triples
\begin{equation}\label{skew sym powers1}
0\to  q_i^*\wedge^j(Sym^dS_i)\overset{\Theta_{ij}}\to
p_i^*\wedge^j(Sym^dS_{i+1})\to\Lambda_{ij}\to0, \ \ \
\Lambda_{ij}:=\coker\Theta_{ij},\ \ \ 1\le j\le \rk T,
\end{equation}
where $\Theta_{ij}$ are the monomorphisms induced by $\theta_i$ in (\ref{taut}). After
restriction to $\mathbb{P}^{n-i}_y$, using (\ref{sym powers1}) and (\ref{sym powers2})
we obtain
\begin{equation}\label{skew sym powers4}
\Lambda_{ij}|_{\mathbb{P}^{n-i}_y}\simeq
\overset{u_{ij}}{\underset{q=1}\oplus}\mathcal{O}_{\mathbb{P}^{n-i}_y}(b_q),\ \ \
-jd\le b_q\le-1,\ \ \ 1\le q\le u_{ij},
\end{equation}
where $u_{ij}:=\binom{s_i+t_i}{j}-\binom{t_i}{j}$. The key observation is that
(\ref{k(n-k)>...}) and (\ref{skew sym powers4}) imply that
$H^a(\Lambda_{ij}|_{\mathbb{P}^{n-i}_y})=0,\ a\ge0,\ 1\le j\le \rk T,\ 1\le i\le k$.
This shows that the Leray spectral sequence
$E_2^{aa'}=H^a(R^{a'}q_{i*}\Lambda_{ij})\Rightarrow H^{\cdot}(\Lambda_{ij})$
degenerates and thus gives
\begin{equation}\label{Leray vanish}
H^a(\Lambda_{ij})=0,\ a\ge0,\ 1\le j\le \rk T,\ 1\le i\le k.
\end{equation}
Since (it is well known that) $H^a(\wedge^j(Sym^dS_i))=H^a(q_i^*\wedge^j(Sym^dS_i)),\
H^a(\wedge^j(Sym^dS_{i+1}))=H^a(p_i^*\wedge^j(Sym^dS_{i+1})),\ a\ge0$, we derive from
(\ref{Leray vanish}) and (\ref{skew sym powers1}) that
\begin{equation}\label{Leray vanish2}
H^a(\wedge^j(Sym^dS_{i+1}))=H^a(\wedge^j(Sym^dS_i)),\ \ \ 1\le i\le k.
\end{equation}
Moreover, setting $j_i:=\rk Sym^dS_i$, we obtain
$\wedge^{j_i}(Sym^dS_i)\simeq\mathcal{O}_{G(i,V)}(-\binom{d+i-1}{i})$, so that,
similarly to (\ref{Leray vanish}), $H^a(\wedge^{j_i}(Sym^dS_i))=0,\ \ a\ge0, \ 1\le
i\le k$. This together with (\ref{Leray vanish2}) yields (\ref{vanishHj}).

Now (\ref{resolutionOBX}) and (\ref{vanishHj}) show that
$h^0(\mathcal{O}_{B_H})=h^0(\mathcal{O}_{G(k+1,V)})=1$. Hence, $B_H$ is connected.
This together with the smoothness of $B_H$ yields its irreducibility.
\end{proof}


Consider the graph of incidence $\Sigma_H=\{(x,\mathbb{P}^k)\in H\times B_H\ |\
x\in\mathbb{P}^k\}$ with projections
\begin{equation}\label{pi1,pi2}
H\overset{\pi_1}\leftarrow\Sigma_H\overset{\pi_2}\rightarrow B_H.
\end{equation}
Since the fibers of $\pi_2$ are isomorphic to $\mathbb{P}^k$, the irreducibility of
$B_H$ implies the irreducibility of $\Sigma_H$.
\begin{lemma}\label{filled}
Let $H\subset\mathbb{P}^n$ be a smooth hypersurface which is generic in the sense of
(\ref{SpanPhi=}). Then

(i) $H$ is filled by the subspaces $\mathbb{P}^k$ of the family $B_H$, and for an
arbitrary $x\in H$ the set $B_H(x):=\pi_2(\pi_1^{-1}(x))$ is equidimensional of
dimension
\begin{equation}\label{dim BX(x)}
\dim B_H(x)=k(n-k)-\binom{k+d}{d}+1;
\end{equation}
moreover, for a generic $x\in H$ $B_H(x)$ is an irreducible subvariety of $B_H$;

(ii) the subset $K_{H,k}(x):=\pi_1(\pi_2^{-1}(B_H(x)))= \underset{\mathbb{P}^k\in
B_H(x)}\cup\mathbb{P}^k$ of $H$ has dimension
\begin{equation}\label{dim KX(x)}
\dim K_{H,k}(x)\ge n-d;
\end{equation}
moreover, for a generic $x\in H$ $K_{H,k}(x)$ is an irreducible subvariety of $H$.
\end{lemma}
\begin{proof}
(i) Let $(x,\mathbb{P}^k)\in\Sigma_H$. Consider the standard Koszul resolution of the
ideal sheaf $\mathcal{I}_{x,\mathbb{P}^k}$
\begin{equation}\label{Koszul0}
0\to\mathcal{O}_{\mathbb{P}^k}(-k)\to...\to
\mathcal{O}_{\mathbb{P}^k}(-i)^{\binom{k}{i}}\to... \to
\mathcal{O}_{\mathbb{P}^k}(-1)^k\to\mathcal{I}_{x,\mathbb{P}^k}\to0.
\end{equation}
Twisting (\ref{Koszul0}) by $N_{\mathbb{P}^k/H}$, we obtain the exact sequence
\begin{equation}\label{Koszul}
0\to N_{\mathbb{P}^k/H}(-k)\to...\to N_{\mathbb{P}^k/H}(-i)^{\binom{k}{i}}\to... \to
N_{\mathbb{P}^k/H}(-1)^k\to \mathcal{I}_{x,\mathbb{P}^k}\otimes
N_{\mathbb{P}^k/H}\to0.
\end{equation}
Since $h^i(N_{\mathbb{P}^k/H}(-i))=0,\ \ 1\le i\le k$ (for $i>1$ this follows
immediately from (\ref{normal bdls}); for $i=1$ see (\ref{h0(N)})), (\ref{Koszul})
gives
\begin{equation}\label{h1(Ix)}
h^1(\mathcal{I}_{x,\mathbb{P}^k}\otimes N_{\mathbb{P}^k/H})=0.
\end{equation}

Next, consider the exact triple
\begin{equation}\label{triple(Ix)}
0\to\mathcal{I}_{x,\mathbb{P}^k}\otimes N_{\mathbb{P}^k/H}\to N_{\mathbb{P}^k/H}\to
\mathbb{C}_x\otimes N_{\mathbb{P}^k/H}\to0.
\end{equation}
Since $\mathbb{C}_x\otimes N_{\mathbb{P}^k/H}\simeq\mathbb{C}^{n-1-k}$, it follows
from (\ref{h0(N)}), (\ref{triple(Ix)}) and (\ref{h1(Ix)}) that
\begin{equation}\label{h0(Ix)}
h^0(\mathcal{I}_{x,\mathbb{P}^k}\otimes N_{\mathbb{P}^k/H})=k(n-k)-\binom{k+d}{d}+1.
\end{equation}
Note that $H^0(\mathcal{I}_{x,\mathbb{P}^k}\otimes N_{\mathbb{P}^k/H})$ is the Zariski
tangent space to $B_H(x)$ at the point $\mathbb{P}^k$. Moreover, (\ref{h1(Ix)}) and
(\ref{h0(Ix)}) imply via deformation theory the smoothness of $B_H(x)$ at
$\mathbb{P}^k$ and the equidimensionality of $B_H(x)$ together with the equality
(\ref{dim BX(x)}). This latter equality shows that $\dim\pi_1(\Sigma_H)=\dim H$. Since
$H$ is irreducible, $\pi_1:\Sigma_H\to H$ is surjective as it is a projective morphism
of projective varieties. This means that $H$ is filled by the spaces $\mathbb{P}^k\in
B_H$.

(ii) Now let $y$ be an arbitrary point of $\mathbb{P}^k$ distinct from $x$ and let
$\mathbb{P}^1$ be a projective line in $\mathbb{P}^k$ joining the points $x$ and $y$.
Twisting (\ref{normal bdls}) by the sheaves $\mathcal{I}_{\mathbb{P}^1,\mathbb{P}^k}$
and $\mathcal{O}_{\mathbb{P}^1}(-2)$ yields the exact triples
\begin{equation}\label{normal bdls1}
0\to\mathcal{I}_{\mathbb{P}^1,\mathbb{P}^k}\otimes N_{\mathbb{P}^k/H}
\to\mathcal{I}_{\mathbb{P}^1,\mathbb{P}^k}\otimes\mathcal{O}_{\mathbb{P}^k}(1)^{n-k}
\overset{\epsilon_k}\to
\mathcal{I}_{\mathbb{P}^1,\mathbb{P}^k}\otimes\mathcal{O}_{\mathbb{P}^k}(d)\to0,
\end{equation}
\begin{equation}\label{normal bdls2}
0\to N_{\mathbb{P}^k/H}\otimes\mathcal{O}_{\mathbb{P}^1}(-2)\to
\mathcal{O}_{\mathbb{P}^1}(-1)^{n-k}\to\mathcal{O}_{\mathbb{P}^1}(d-2)\to0.
\end{equation}

Consider the morphism $\epsilon_k$ in (\ref{normal bdls1}). Passing to sections, we
obtain the homomorphism
$H^0(\epsilon_k):H^0(\mathcal{I}_{\mathbb{P}^1,\mathbb{P}^k}\otimes
\mathcal{O}_{\mathbb{P}^k}(1)^{n-k})\to
H^0(\mathcal{I}_{\mathbb{P}^1,\mathbb{P}^k}\otimes \mathcal{O}_{\mathbb{P}^k}(d))$. To
show that $H^0(\epsilon_k)$ is an epimorphism, consider the standard Koszul resolution
of the sheaf
$\mathcal{I}_{\mathbb{P}^1,\mathbb{P}^k}\otimes\mathcal{O}_{\mathbb{P}^k}(1)$
$$
0\to\mathcal{O}_{\mathbb{P}^k}(2-k)\to...\to\mathcal{O}_{\mathbb{P}^k}(-1)^{\binom{k-1}{2}}\to
\mathcal{O}_{\mathbb{P}^k}^{k-1}\overset{e_1}\to
\mathcal{I}_{\mathbb{P}^1,\mathbb{P}^k}\otimes\mathcal{O}_{\mathbb{P}^k}(1)\to0.
$$
Passing to cohomology, we obtain the epimorphism in sections
$H^0(e_1):H^0(\mathcal{O}_{\mathbb{P}^k}^{k-1})\to
H^0(\mathcal{I}_{\mathbb{P}^1,\mathbb{P}^k}\otimes\mathcal{O}_{\mathbb{P}^k}(1))$.
Twisting the above resolution by $\mathcal{O}_{\mathbb{P}^k}(d-1)$ and again passing
to cohomology, we obtain an epimorphism
$H^0(e_d):H^0(\mathcal{O}_{\mathbb{P}^k}(d-1)^{k-1})\to
H^0(\mathcal{I}_{\mathbb{P}^1,\mathbb{P}^k}\otimes\mathcal{O}_{\mathbb{P}^k}(d))$. Now
the homomorphisms $H^0(\epsilon_k)$, $H^0(e_1)$ and $H^0(e_d)$ fit in a commutative
diagram
$$
\xymatrix{ H^0(\mathcal{I}_{\mathbb{P}^1,\mathbb{P}^k}\otimes
\mathcal{O}_{\mathbb{P}^k}(1)^{n-k})\ar[r]^{\ \ H^0(\epsilon_k)} &
H^0(\mathcal{I}_{\mathbb{P}^1,\mathbb{P}^k}\otimes\mathcal{O}_{\mathbb{P}^k}(d))\\
H^0(\mathcal{O}_{\mathbb{P}^k}^{(k-1)(n-k)})\ar@{->>}[r]^{H^0(\epsilon_k)^{k-1}}
\ar@{->>}[u]^{H^0(e_1)^{n-k}} &
H^0(\mathcal{O}_{\mathbb{P}^k}(d-1)^{k-1}),\ar@{->>}[u]_{H^0(e_d)} }
$$
in which the surjectivity of the lower horizontal map $H^0(\epsilon_k)^{k-1}$ follows
from (\ref{SpanPhi=}). Hence $H^0(\epsilon_k)$ is an epimorphism. Thus the cohomology
sequence of (\ref{normal bdls1}) yields
\begin{equation}\label{h^0(JN)}
h^0(\mathcal{I}_{\mathbb{P}^1,\mathbb{P}^k}\otimes N_{\mathbb{P}^k/H})=
(k-1)(n-k)-\binom{k+d}{d}+d+1,\ \ \ h^1(\mathcal{I}_{\mathbb{P}^1,\mathbb{P}^k}\otimes
N_{\mathbb{P}^k/H})=0.
\end{equation}
Next, (\ref{normal bdls2}) implies
$h^0(N_{\mathbb{P}^k/H}\otimes\mathcal{O}_{\mathbb{P}^1}(-2))=0, \
h^1(N_{\mathbb{P}^k/H}\otimes\mathcal{O}_{\mathbb{P}^1}(-2))=d-1$. This together with
(\ref{h^0(JN)}) and the exact triple
\begin{equation}\label{exactP1}
0\to\mathcal{I}_{\mathbb{P}^1,\mathbb{P}^k}\otimes N_{\mathbb{P}^k/H}
\to\mathcal{I}_{x\cup y,\mathbb{P}^k}\otimes N_{\mathbb{P}^k/H}
\to\mathcal{O}_{\mathbb{P}^1}(-1)\otimes N_{\mathbb{P}^k/H}\to0.
\end{equation}
yields
\begin{equation}\label{hi(Ixy)}
h^0(\mathcal{I}_{x\cup y,\mathbb{P}^k}\otimes
N_{\mathbb{P}^k/H})=(k-1)(n-k)-\binom{k+d}{d}+d+1, \ \ \ h^1(\mathcal{I}_{x\cup
y,\mathbb{P}^k}\otimes N_{\mathbb{P}^k/H})=d-1.
\end{equation}

Put $\Sigma_H(x):=\pi_2^{-1}(B_H(x)),\ \ \pi_1(x):=\pi_1|_{\Sigma_H(x)},$ and let
$$
K_{H,k}(x)\overset{\pi_1(x)}\leftarrow\Sigma_H(x)\overset{\pi_2(x)}\rightarrow B_H(x)
$$
be the diagram of projections. For any $y\in K_{H,k}(x),\ y\ne x$, consider the fibre
$B_{H,x}(y):=\pi_1(x)^{-1}(y)$ as lying in $B_H(x)$. The Zariski tangent space to
$B_{H,x}(y)$ at the point $\mathbb{P}^k$ coincides with $H^0(\mathcal{I}_{x\cup
y,\mathbb{P}^k}\otimes N_{\mathbb{P}^k/H})$, hence by (\ref{hi(Ixy)}) and deformation
theory we have
\begin{equation}\label{dim BXx(y)}
(k-1)(n-k)-\binom{k+d}{d}+d+1\ge\dim B_{H,x}(y)\ge(k-1)(n-k)-\binom{k+d}{d}+2.
\end{equation}
Clearly $\dim B_{H,x}(y)>0$, hence $\pi_1(x)$ is surjective. Since the fibre of
$\pi_2(x)$ is $\mathbb{P}^k$, this, together with (\ref{dim BX(x)}), (\ref{dim
BXx(y)}) and the irreducibility of $B_H(x)$, implies (\ref{dim KX(x)}) and the
irreducibility of $K_{H,k}(x)$.
\end{proof}

As a corollary of this lemma we obtain the following theorem.
\begin{theorem}\label{any X is filled}
Any hypersurface $H$ of degree $d$ in $\mathbb{P}^n$ is filled by subspaces
$\mathbb{P}^k\subset\mathbb{P}^n$.
\end{theorem}

\begin{proof}
Consider the graphs of incidence $\Pi:=\{(\mathbb{P}^k,x)\in
G(k+1,V)\times\mathbb{P}^n\ |x\in\mathbb{P}^k\}$ and $\tilde{H}:=\{(H,x)\in
\mathbb{P}^s\times\mathbb{P}^n\ |\ x\in H\}$ fitting in the commutative diagram
\begin{equation}\label{graph Pi}
\xymatrix{ \Pi\ar@{_{(}->}[d] & \Pi_\Gamma\ar@{_{(}->}[d]\ar[r]^{pr_2}\ar[l]_{pr_1} &
\tilde{H}\ar@{_{(}->}[d]\\
G(k+1,V)\times\mathbb{P}^n & \Gamma\times\mathbb{P}^n \ar[r]^{\tilde{q}\times id\ \
}\ar[l]_{\ \ \ \ \ \ \tilde{p}\times id} & \mathbb{P}^s\times\mathbb{P}^n,}
\end{equation}
where $\Gamma,\ \tilde{p}$ and $\tilde{q}$ were defined in (\ref{p,q}),
$\Pi_\Gamma=(\tilde{p}\times id)^{-1}(\Pi)$ and $pr_1$ and $pr_2$ are the induced
projections. Since a generic smooth $H\in\mathbb{P}^s$ is filled by projective
subspaces $\mathbb{P}^k\subset\mathbb{P}^n$ (Lemma \ref{filled}(i)), $pr_2$ is
dominant. Hence $pr_2$ is surjective since all varieties and morphisms in (\ref{graph
Pi}) are projective. This implies the statement.
\end{proof}



\end{sub}


\begin{sub}{\bf Projective subspaces in varieties of bounded codimension and degree and of
growing dimension.}\label{varieties} \rm Let $X\subset\mathbb{P}^n$ be an irreducible
projective variety satisfying the conditions
\begin{equation}\label{codimX,degX}
1\le c:=\codim_{\mathbb{P}^n}X, \ \ \ \ \deg X\le d,
\end{equation}
where $c$ is a constant. Assume that $\mathbb{P}^n=\Span X$. Then it is well known
that $\deg X\ge c+1$. If $c\ge2$, starting with $X_0:=X$ one can construct inductively
a sequence of projective varieties $X_i\subset\mathbb{P}^{n-i}, \ 0\le i\le c-1,$ of
respective codimensions $c-i$, together with linear projections
$$
p_{x_i}:\mathbb{P}^{n-i}\dasharrow\mathbb{P}^{n-i-1}, \ \ \ 0\le i\le c-2,
$$
with centers at points $x_i\in X_i\smallsetminus\Sing X_i$ such that each restriction
\begin{equation}\label{bir pi}
p_i:=p_{x_i}|_{X_i}:X_i\dasharrow X_{i+1}, \ \ \ 0\le i\le c-2.
\end{equation}
is a birational isomorphism. For this, it suffices to fix $x_i\in
X_i\smallsetminus\Sing X_i$ and let $X_{i+1}$ be the closure of $p_{x_i}(X_i)$ in
$\mathbb{P}^{n-i}$. Then $\deg X_{i+1}=\deg X_i-1$. The fact that $p_i$ is birational
is standard.


Next, using the notation (\ref{k(n-k)>...}), we set
$$
k_{c-1}(n):=k(n-c+1),\ \ \ k_{c-1-i}(n):=\underset{i}{\underbrace{\Big[\frac{1}{2}...
\Big[\frac{1}{2}\Big[\frac{1}{2}k_{c-1}(n)\Big]\Big]...\Big]}},\ \ \ 1\le i\le c-1.
$$

We now argue by reverse induction that $X=X_0$ is filled by projective subspaces of
dimension $k_0(n)$. By definition, $X_{c-1}$ is a hypersurface in
$\mathbb{P}^{n-(c-1)}$ of degree
\begin{equation}\label{Xc-1}
\deg X_{c-1}=\deg X-(c-1)\le d.
\end{equation}
Hence, by Theorem \ref{any X is filled}, $X_{c-1}$ is filled by subspaces
$\mathbb{P}^{k_{c-1}(n)}$ of $\mathbb{P}^{n-(c-1)}$. This settles the base of
induction.

For the induction step, consider the birational map (\ref{bir pi}). Assume that
$X_{i+1}$ is filled by subspaces $\mathbb{P}^{k_{i+1}(n)}\subset\mathbb{P}^{n-i-1}$.
Let $B$ be an irreducible component of the base of all such subspaces, with the
property that the subspaces in $B$ fill $X_{i+1}$. Take a generic space
$\mathbb{P}^{k_{i+1}(n)}\in B$ and consider the closure
$Y_{i+1}:=\overline{p_i^{-1}(\mathbb{P}^{k_{i+1}(n)})}. $ Since $\mathbb{P}^{k_{i+1}}$
is a generic point of $B$, the rational map
$\tilde{p}:=p_i|_{Y_{i+1}}:Y_{i+1}\dasharrow\mathbb{P}^{k_{i+1}(n)}$ is a linear
projection from the point $x_i\in Y_{i+1}$, and one of the following alternatives
holds.

(i) $Y_{i+1}$ is an irreducible quadric and
$$
\tilde{p}:Y_{i+1}\dasharrow\mathbb{P}^{k_{i+1}(n)}
$$
is a birational (stereographic) projection from the point $x_i\in Y_{i+1}$.

(ii) $Y_{i+1}$ is a reducible quadric containing as a component a certain
$k_{i+1}(n)$-dimensional space $\tilde{\mathbb{P}}^{k_{i+1}(n)}$ mapping
isomorphically onto $\mathbb{P}^{k_{i+1}(n)}$,
$$
\tilde{p}:\tilde{\mathbb{P}}^{k_{i+1}(n)}\overset{\sim}\to\mathbb{P}^{k_{i+1}(n)}.
$$
Consider these two cases.

In case (i) $Y_{i+1}$ is an irreducible quadric of dimension $k_{i+1}(n)$ filled by
projective spaces of dimension $\Big[\frac{1}{2}k_{i+1}(n)\Big]=k_i(n)$. Since
$\mathbb{P}^{k_{i+1}(n)}$ is a generic point of $B$, the quadrics
$Y_{i+1}(\mathbb{P}^{k_{i+1}(n)}),\ \ \mathbb{P}^{k_{i+1}(n)}\in B,$ fill the variety
$X_i$. Hence, the subspaces $\mathbb{P}^{k_i(n)}$ fill $X_i$.

In case (ii) the irreducibility of $X_i$ and the birationality of $p_i$ imply that the
subspaces $\tilde{\mathbb{P}}^{k_{i+1}(n)}$ fill $X_i$. Moreover, each
$\tilde{\mathbb{P}}^{k_{i+1}(n)}$ is filled by subspaces $\mathbb{P}^{k_i(n)}$. Hence
$X_i$ is filled by these $\mathbb{P}^{k_i(n)}$'s as well.

Finally note that  $\underset{n\to\infty}\lim k_0(n)=\infty$. We have thus proved the
following theorem.

\begin{theorem}\label{X is filled II}
Let $X\subset\mathbb{P}^n$ be an irreducible projective variety satisfying the
conditions (\ref{codimX,degX}) and ${\rm Span} X=\mathbb{P}^n$. Then $X$ is filled by
projective subspaces $\mathbb{P}^{k_0(n)}\subset\mathbb{P}^n$ with
$\underset{n\to\infty}\lim k_0(n)=\infty$.
\end{theorem}

\end{sub}



\begin{sub}{\bf Chains of projective subspaces connecting the points of varieties of bounded
codimension and degree.}\label{chains} \rm Let again $H$ be a smooth hypersurface of
degree $d\ge2$ in $\mathbb{P}^n$ and $x\in H$. Denote by $\mathbb{P}^{n-1}(x)$ the
hyperplane in $\mathbb{P}^n$ tangent to $H$ at the point $x$. Take an affine subspace
$\mathbb{A}^{n-1}(x)$ of $\mathbb{P}^{n-1}(x)$ containing $x$, together with affine
coordinates $(y_1,...,y_{n-1})$ around $x$ in $\mathbb{A}^{n-1}(x)$. The intersection
$Y_H(x):=H\cap\mathbb{A}^{n-1}(x)$ is a hypersurface in $\mathbb{A}^{n-1}(x)$ and is
given by an equation $\Psi_x=0$ for some polynomial $\Psi_x=\Psi_x(y_1,...,y_{n-1})$
of degree $d$. Decompose $\Psi_x$ into a sum of its homogeneous components
\begin{equation}\label{homog}
\Psi_x={\underset{p=2}{\overset{d}\sum}}\Psi_p(y_1,...,y_{n-1}),\ \ \ \deg\Psi_p=p.
\end{equation}
Consider $(y_1:y_2:...:y_{n-1})$ as homogeneous coordinates in $\mathbb{P}^{n-2}$;
respectively, consider $\Psi_p$ as forms $\Psi_p\in
H^0(\mathcal{O}_{\mathbb{P}^{n-2}}(p))$. Define the closed subset
$$
X_x={\underset{p=2}{\overset{d}\bigcap}}\{\Psi_p(y_1,...,y_{n-1})=0\},\ \ \
\deg\Psi_p=p.
$$
in $\mathbb{P}^{n-2}$. Then Bezout's Theorem implies
\begin{equation}\label{codimZ, degZ}
\codim_{\mathbb{P}^{n-2}}X\le d-1,\ \ \ \deg X\le d!
\end{equation}
for any irreducible component $X$ of $X_x$. Therefore $n-2\ge\dim\Span X\ge n-d-1$. In
particular, the codimension and degree of $X$ are bounded by constants not depending
on $n$, hence Theorem \ref{X is filled II} applies to $X$. This proves the following
lemma.
\begin{lemma}\label{Z}
There exists $n(d)\in\mathbb{Z}_{>0}$ such that, for $n\ge n(d)$, the variety $X_x$ is
connected and any irreducible component $X$ of $X_x$ is filled by subspaces
$\mathbb{P}^{\tilde{k}(n)}\subset\mathbb{P}^{n-2}$ with
$\underset{n\to\infty}\lim\tilde{k}(n)=\infty$.
\end{lemma}
Let $K_H(x)$ be the cone in $\mathbb{A}^{n-1}(x)$ over $X_x$. By Lemma \ref{Z} the
closure $\overline{K_H(x)}$ of $K_H(x)$ in $\mathbb{P}^{n-1}(x)$ is filled by
subspaces $\mathbb{P}^{\tilde{k}(n)}$. Now consider the subvariety $K_{H,k}(x)$ of $H$
defined in Lemma \ref{filled}(ii). By definition, $K_{H,k}(x)$ is filled by those
subspaces $\mathbb{P}^k$ filling $H$ which pass through $x$. Clearly,
\begin{equation}\label{KXk(x) in KX(x)}
K_H(x)\supset K_{H,k}(x)\ \ \ {\rm for}\ \ \ k=\tilde{k}(n).
\end{equation}

Assume that $H\subset\mathbb{P}^n$ is a generic smooth hypersurface of degree $d$ and
$x$ is a generic point of $H$. In particular, the forms $\Psi_p$ are generic points of
the spaces $H^0(\mathcal{O}_{\mathbb{P}^{n-2}}(p))$. Hence, in this case $X=X_x$ is a
smooth and irreducible complete intersection of $d-1$ hypersurfaces $\{\Psi_p=0\},\
p=2,...,d$ in $\mathbb{P}^{n-2}$, and the inequalities (\ref{codimZ, degZ}) become
equalities. This together with (\ref{KXk(x) in KX(x)}) and (\ref{dim KX(x)}) implies
$K_H(x)=K_{H,k}(x)$. In addition, the sheaf $\mathcal{O}_{X}$ has a standard Koszul
resolution $0\to\mathcal{O}_{\mathbb{P}^{n-1}}(2-d(d+1)/2)\to...\to
\overset{d}{\underset{p=2}\oplus}\mathcal{O}_{\mathbb{P}^{n-2}}(1-p)
\to\mathcal{O}_{\mathbb{P}^{n-2}}(1)\overset{res}\to\mathcal{O}_{X}(1)\to0.$ This
resolution together with (\ref{k(n-k)>...}) shows that the restriction map
$H^0(res):H^0(\mathcal{O}_{\mathbb{P}^{n-2}}(1))\to H^0(\mathcal{O}_{X}(1))$ is an
isomorphism. Therefore $\Span X=\mathbb{P}^{n-2}$ and, consequently,
\begin{equation}\label{Span K}
\Span K_{H,k}(x)=\mathbb{P}^{n-1}(x).
\end{equation}

We now define a sequence of irreducible subvarieties $x\in X_1\subset
X_2\subset...\subset X_i\subset...\subset H$ by induction:

1) $X_1:=K_{H,k}(x)$;

2) $X_{i+1}:=\pi_1(\pi_2^{-1}(Y_i))$ for $i\ge1$,  $Y_i$ being any irreducible
component of $\pi_2(\pi_1^{-1}(X_i))$, where $\pi_1$ and $\pi_2$ are introduced in
diagram (\ref{pi1,pi2}).

Since $X$ is irreducible, this sequence stabilizes, i.e.
\begin{equation}\label{again Xi's}
X_1\subset X_2\subset...\subset X_{i_0}=X_{i_0+1}...\subset H.
\end{equation}
for some $i_0$. Consider the dense open subset $U:=\{x'\in H\ |\ K_{H,k}(x')$ is
irreducible and $\Span(K_{H,k}(x'))=\mathbb{P}^{n-1}(x')\}\subset H$. By construction,
$x\in U$, hence $X_{i_0}\cap U$ is a dense open subset of $X_{i_0}$. Moreover, by the
definition of $X_{i_0}$  we have
\begin{equation}\label{KX,k(x')}
K_{H,k}(x')\subset X_{i_0}
\end{equation}
for $x'\in X_{i_0}\cap U$. Denote by $H(x')$ the projective subspace of $\mathbb{P}^n$
tangent to $X_{i_0}$ at the point $x'\in(X_{i_0}\smallsetminus\Sing X_{i_0})\cap U$.
Since $K_{H,k}(x')$ is by definition filled by projective subspaces on $H$ through
$x'$, it follows from (\ref{KX,k(x')}) that $K_{H,k}(x')\subset
H(x')\subset\mathbb{P}^{n-1}(x')$. On the other hand, since $x'\in U$, it follows that
$\Span K_{H,k}(x)=H(x')$. As $H(x')$ is a subspace of $\mathbb{P}^{n-1}(x')$, by
(\ref{Span K}) we have $H(x')=\mathbb{P}^{n-1}(x')$. Hence, since $x'$ is a
nonsingular point of $X_{i_0}$, we obtain $\dim X_{i_0}=\dim H$, so that
\begin{equation}\label{Xi0=X}
X_{i_0}=H.
\end{equation}
This equality and the construction of the chain (\ref{again Xi's}) shows that the
point $x\in H$ can be joined with any point $x'\in H$ by a chain of subspaces
$\mathbb{P}^k_1,\mathbb{P}^k_2,...,\mathbb{P}^k_{i_0}$. We thus have
\begin{equation}\label{join xx'}
x\in\mathbb{P}^k_1\subset\mathbb{P}^k_1\cup\mathbb{P}^k_2\cup...\cup\mathbb{P}^k_{i_0}
\supset\mathbb{P}^k_{i_0}\ni x'.
\end{equation}

Finally, we will show that (\ref{join xx'}) holds also without the genericness
assumption on $H$ and $x$. This is done by essentially the same argument as in the
proof of Theorem \ref{any X is filled}. Indeed, consider the Grassmannian
$G:=G(k+1,V)$, the incidence variety
$$
Inc^{i_0}(G):=\{(\mathbb{P}^k_1,...,\mathbb{P}^k_{i_0})\in G^{\times i_0}\ |\
\mathbb{P}^k_1,...,\mathbb{P}^k_{i_0}\ {\rm is\ a\ chain\ of\ subspaces\ of}\
\mathbb{P}^n\},
$$
and the graphs of incidence
$$
\Pi_{i_0}:=\{(\mathbb{P}^k_1,...,\mathbb{P}^k_{i_0},x,x')\in
Inc^{i_0}(G)\times\mathbb{P}^n\times\mathbb{P}^n\ | x\in\mathbb{P}^k_1,\
x'\in\mathbb{P}^k_{i_0}\},
$$
$$
\widetilde{\widetilde{H}}:=\{(H,x,x')\in
\mathbb{P}^s\times\mathbb{P}^n\times\mathbb{P}^n\ | \ x,x'\in H\},
$$
$$
\Gamma_{i_0}:=\{(\mathbb{P}^k_1,...,\mathbb{P}^k_{i_0},G)\in
Inc^{i_0}(G)\times\mathbb{P}^s\ |\ \mathbb{P}^k_1,...,\mathbb{P}^k_{i_0}\subset H\}
$$
with natural projections
$$
Inc^{i_0}(G)\overset{\tilde{p}_{i_0}}\leftarrow\Gamma_{i_0}
\overset{\tilde{q}_{i_0}}\to\mathbb{P}^s.
$$
We have the commutative diagram
\[\xymatrix{
\Pi_{i_0}\ar@{_{(}->}[d] & \Pi_{\Gamma_{i_0}}\ar@{_{(}->}[d]\ar[r]^{pr_2}\ar[l]_{pr_1}
&
\widetilde{\widetilde{H}}\ar@{_{(}->}[d]\\
Inc^{i_0}(G)\times\mathbb{P}^n\times\mathbb{P}^n &
\Gamma_{i_0}\times\mathbb{P}^n\times\mathbb{P}^n \ar[r]^{\tilde{q}_{i_0}\times id\ \
}\ar[l]_{\ \ \ \ \ \ \tilde{p}_{i_0}\times id} &
\mathbb{P}^s\times\mathbb{P}^n\times\mathbb{P}^n,}
\]
where $pr_1$ and $pr_2$ are the induced projections. As a generic smooth
$H\in\mathbb{P}^s$ is filled by projective subspaces $\mathbb{P}^k\subset\mathbb{P}^n$
(Lemma \ref{filled}(i)), the morphism $pr_2$ is dominant. Hence $pr_2$ is surjective
since all varieties and morphisms in the above diagram are projective. This is
equivalent to ( \ref{join xx'}) for any $H\in\mathbb{P}^s$ and any $x,x'\in H$.

We thus have proved the following lemma.

\begin{lemma}\label{join xx' II}
Let $H$ be a hypersurface of degree $d$ in $\mathbb{P}^n$. Any two distinct points
$x,x'\in H$ can be joined by a chain (\ref{join xx'}) of subspaces $\mathbb{P}^k$ of
$H$.
\end{lemma}

Finally, Lemma \ref{join xx' II} together with Theorem \ref{X is filled II} leads to
our main result in section 3.
\begin{theorem}\label{join xx' III}
Under the assumptions of Theorem \ref{X is filled II}, any two distinct points
$x,x'\in X$ can be joined by a chain (\ref{join xx'}) of subspaces
$\mathbb{P}^{k_0(n)}$ of $X$ with $\underset{n\to\infty}\lim k_0(n)=\infty$.
\end{theorem}

\end{sub}



\vspace{1cm}

\section{A sufficient condition on $m$ for a vector bundle on $\mathbb{P}^N$ to be $m$-regular}
\label{m-reg} \vspace{0.5cm}

Recall that a vector bundle $E$ on a scheme $Y$ is called {\it ample} if the
invertible Grothendieck sheaf $\mathcal{O}_{\mathbb{P}(E^\vee)}(1)$ on
$\mathbb{P}(E^\vee)$ is ample. The following result is well known - see, e.g.,
\cite[Prop. 6.3.56]{L}.
\begin{lemma}\label{ample}
Let $E$ be a vector bundle on $\mathbb{P}^N$. Then $E(a)$ is ample for any
$a\in\mathbb{Z}_{\ge a_0}$, $a_0$ being some fixed integer.
\end{lemma}

\begin{lemma}\label{ample on P1}
Let $E$ be a vector bundle on $\mathbb{P}^1$. Then $E(a)$ is ample for
$a\ge1-\delta_B(E)$.
\end{lemma}
\begin{proof}
By Grothendieck's theorem,
$E\simeq\overset{r}{\underset{j=1}\oplus}\mathcal{O}_{\mathbb{P}^1}(a_j)$, where
$\delta_B(E)=a_1\le a_2\le...\le a_r$, $r=\rk E$. Hence, for $a\ge1-\delta_B(E)$, the
bundle $E(a)$ is a direct sum of ample line bundles. By \cite [Prop. 6.1.12(i)]{L}
$E(a)$ is itself ample.
\end{proof}

We now recall the notion of {\it degree} of a vector bundle $\mathcal{E}$ on a
1-dimensional scheme $Y$. If $Y$ is a smooth irreducible curve,
$\deg\mathcal{E}:=\chi(\mathcal{E})-\chi(\mathcal{O}_Y)\rk\mathcal{E}$. If $Y$ is
irreducible, but not necessarily smooth, the degree $\deg\mathcal{E}$ is defined as
the degree of the pullback of $\mathcal{E}$ to the normalization of $Y$. If $Y$ is a
general 1-dimensional scheme with irreducible components $Y_1,...,Y_q$, then the
multiplicities $k_i\in\mathbb{Z}_{>0}$ of $Y_i$ in $Y$ are well defined (see
\cite[1.5]{F}), and we set
\begin{equation}\label{degE}
\deg\mathcal{E}=\underset{i}\sum k_i\deg(\mathcal{E}|_{Y_i}).
\end{equation}

\begin{lemma}\label{deg>0}
Let $E$ be a vector bundle on $\mathbb{P}^N$ and let
$pr:\mathbb{P}(E^\vee)\to\mathbb{P}^N$ be the projection. Let $Y$ be a 1-dimensional
subscheme of $\mathbb{P}(E^\vee)$ such that $Y_{red}\subset pr^{-1}(\mathbb{P}^1)$ for
some line $\mathbb{P}^1\subset\mathbb{P}^N.$ Consider the line bundle
$L_0=\mathcal{O}_{\mathbb{P}(E^\vee)}(1)\otimes pr^*\mathcal{O}_{\mathbb{P}^1}(a)$ on
$\mathbb{P}(E^\vee)$ for $a\ge1-\delta_B(E)$. Then
\begin{equation}\label{L0|Y}
\deg(L_0|_Y)>0.
\end{equation}
\end{lemma}
\begin{proof}
By (\ref{degE}),
\begin{equation}\label{again L0|Y}
\deg(L_0|_Y)=\underset{i}\sum k_i\deg(L_0|_{Y_i}),\ \ k_i>0,
\end{equation}
where $Y_i$ are the irreducible components of $Y$. Since
$\delta_B(E|_{\mathbb{P}^1})\ge \delta_B(E)$, it follows from Lemma \ref{ample on P1}
that the sheaf $L_0|_{pr^{-1}(\mathbb{P}^1)}$ is ample. Hence $\deg(L_0|_{Y_i})>0$ for
each $Y_i$ above, and (\ref{again L0|Y}) implies (\ref{L0|Y}).
\end{proof}

Let $Z_1$ be an arbitrary reduced irreducible curve in $\mathbb{P}^N$ with $N\ge3$.
Pick a projective line $l_0\subset\mathbb{P}^N$ and a subspace
$\mathbb{P}^{N-2}\subset\mathbb{P}^N$ such that
\begin{equation}\label{empty}
l_0\cap Z_1=\mathbb{P}^{N-2}\cap Z_1=\emptyset.
\end{equation}
Fix homogeneous coordinates $(x_0:...:x_N)$ in $\mathbb{P}^N$ so that
$l_0=\{x_2=...=x_N=0\},\ \mathbb{P}^{N-2}=\{x_0=x_1=0\}$, and fix the isomorphism
$$
\Lambda:\mathbb{C}^*\times\mathbb{P}^N\overset{\sim}\to\mathbb{C}^*\times\mathbb{P}^N,
(t,(x_0:...:x_N))\mapsto(t,(x_0:x_1:tx_2:...:tx_N)),
$$
where $\mathbb{C}^*=\mathbb{C}\smallsetminus\{0\}$. Set
$\Gamma^*:=\Lambda(\mathbb{C}^*\times Z_1)$ and consider the Hilbert scheme
$\mathcal{H}:={\rm Hilb}^{P_{Z_1}}(\mathbb{P}^N)$, where $P_{Z_1}$ is the Hilbert
polynomial $P_{Z_1}(n)=\chi(\mathcal{O}_{\mathbb{P}^N}(n)|_{Z_1})$. By construction
$\Gamma^*\to\mathbb{C}^*$ is a flat family of curves over $\mathbb{C}^*$, hence it
defines a morphism $g:\mathbb{C}^*\to\mathcal{H}$ such that $\Gamma^*=
\Gamma_\mathcal{H}\times_\mathcal{H}\mathbb{C}^*$, where
$\Gamma_\mathcal{H}\subset\mathbb{P}^N\times \mathcal{H}$ is the universal family of
curves. The coordinate $t$ on $\mathbb{C}^*$ identifies $\mathbb{C}^*$ with
$\mathbb{P}^1\smallsetminus\{z_0,z_\infty\}$, where $z_0=\{t=0\},\ \
z_\infty=\{t=\infty\},$ and since the Hilbert scheme $\mathcal{H}$ is projective, the
morphism $g$ extends to a morphism $\tilde{g}:\mathbb{P}^1\to H$. We thus obtain a
flat family
$\phi:\Gamma=\Gamma_\mathcal{H}\times_\mathcal{H}\mathbb{P}^1\to\mathbb{P}^1$ of
curves over $\mathbb{P}^1$ such that $Z_1=\phi^{-1}(z_1)$ for $z_1:=\{t=1\}$, and
$(\phi^{-1}(z_0))_{red}=l_0$.


Let again $E$ be a vector bundle of rank $\rk E\ge2$ on $\mathbb{P}^N$ and let
$pr:\mathbb{P}(E^\vee)\to\mathbb{P}^N$ be the projection. Consider the projection
$q:\Gamma\to\mathbb{P}^N$ and the scheme
$\Gamma^E:=\mathbb{P}(q^*E^\vee)=\mathbb{P}(E^\vee)\times_{\mathbb{P}^N}\Gamma$ with
projections $\mathbb{P}(E^\vee)\overset{q'}\leftarrow \Gamma^E\overset{pr'}\to\Gamma$
and $\rho=\phi\circ pr':\Gamma^E\to\mathbb{P}^1.$ Note that, by Lemma \ref{ample}
there exists $a_0\in\mathbb{Z}$ such that the line bundle
$A=\mathcal{O}_{\mathbb{P}(E^\vee)}(1)\otimes pr^*O_{\mathbb{P}^N}(a_0)$ is ample on
$\mathbb{P}(E^\vee)$; hence the line bundle $q'^*A$ is $\rho$-ample on $\Gamma^E$.


Fix an irreducible curve $Y_1$ in $\Gamma^E$ such that $pr'(Y_1)=Z_1$, and denote by
$P_{Y_1}$ the Hilbert polynomial $P_{Y_1}(n):=\chi(q'^*A^{\otimes n}|_{Y_1})$.
Consider the relative Hilbert scheme $\mathcal{H}_{\mathbb{P}^1}={\rm
Hilb}^{P_{Y_1}}(\Gamma^E/\mathbb{P}^1)$, together with the natural surjective
projective morphism $f:\mathcal{H}_{\mathbb{P}^1}\to\mathbb{P}^1$ and the universal
family $\Sigma\hookrightarrow \Gamma^E\times_{\mathbb{P}^1}\mathcal{H}_{\mathbb{P}^1}$
with projections
$\Gamma^E\overset{p''}\leftarrow\Sigma\overset{q''}\to\mathcal{H}_{\mathbb{P}^1}$. By
definition, there is a point $y_1\in\mathcal{H}_{\mathbb{P}^1}$ such that
\begin{equation}\label{y1}
q''^{-1}(y_1)\overset{p''}{\underset{\sim}\longrightarrow}Y_1
\end{equation}
and  $f(y_1)=a_1$. Next, consider the normalization $\nu:Z\to Z_1$ of $Z_1$ and the
surfaces $\mathcal{S}=\mathbb{P}(\nu^*(E^\vee|_{Z_1}))$ and
$\mathcal{S}_1=\mathbb{P}(E^\vee|_{Z_1})\subset X_\Gamma$ with their projections
$pr_\mathcal{S}:\mathcal{S}\to Z$ and $pr_{\mathcal{S}_1}:\mathcal{S}_1\to Z_1$. By
construction, the morphism $\nu$ lifts to the normalization
$\tilde{\nu}:\mathcal{S}\to \mathcal{S}_1$ such that
$pr_{\mathcal{S}_1}\circ\tilde{\nu}=\nu\circ pr_\mathcal{S}$, and the curve
$Y=\tilde{\nu}^{-1}(Y_1)$ is a multisection of the projection $pr_\mathcal{S}$.

Consider the Hilbert polynomial $P_Y(n):=\chi(\tilde{\nu}^*q'^*A^{\otimes n}|_{Y_1})$.
Since $\mathcal{S}$ is a smooth surface, the Hilbert scheme ${\rm
Hilb}^{P_Y}(\mathcal{S})$ coincides with the linear series
$|\mathcal{O}_\mathcal{S}(Y)|\simeq\mathbb{P}^h,\ \
h=h^0(\mathcal{O}_\mathcal{S}(Y))-1$, and there is a bijective morphism
$\mathbb{P}^h={\rm Hilb}^{P_Y}(\mathcal{S})\to{\rm
Hilb}^{P_{Y_1}}(\mathcal{S}_1)=f^{-1}(a_1): C\mapsto\tilde{\nu}(C)$. Thus the fibre
$f^{-1}(a_1)$ is irreducible.

Since the morphism $f:\mathcal{H}_{\mathbb{P}^1}\to\mathbb{P}^1$ is projective, the
scheme $\mathcal{H}_{\mathbb{P}^1}$ is projective as well. Therefore, in view of the
surjectivity and flatness of $f$ and the irreducibility of the fibre $f^{-1}(a_1)$,
there exists a smooth irreducible curve $T$ and a morphism
$\theta:T\to\mathcal{H}_{\mathbb{P}^1}$ such that
$\theta_T=f\circ\theta:T\to\mathbb{P}^1$ is surjective. Hence
\begin{equation}\label{t0}
\theta_T(t_0)=z_0
\end{equation}
for some $t_0\in T$, and, since $f(y_1)=z_1$,
\begin{equation}\label{t1}
\theta(t_1)=y_1,\ \ \ \theta_T(t_1)=z_1
\end{equation}
for some $t_1\in T$.

Consider the fibre product $\Sigma_T=\Sigma\times_{\mathbb{P}^1}T$ with projections
$p_T:\Sigma_T\to T$, $q_T:\Sigma_T\to\Sigma\overset{p''}\to X_\Gamma\overset{q''}\to
\mathbb{P}(E^\vee)$, and the embedding $i=(q_T,p_T):\Sigma\hookrightarrow
\mathbb{P}(E^\vee)\times T$. The family $p_T:\Sigma_T\to T$ is a flat family of curves
in $\mathbb{P}(E^\vee)$ with base $T$ such that the fibre $p_T^{-1}(t_1)$ coincides
with $Y_1$, and the reduced fibre $(Y_0)_{red}:=(p_T^{-1}(t_0))_{red}$ lies in
$pr^{-1}(l_0)$. Next, consider the line bundle $L_T=i^*(L_0\boxtimes\mathcal{O}_T)$ on
$\Sigma_T$, where $L_0$ is the line bundle on $X$ defined in Lemma \ref{deg>0}. The
degree $\deg(L_T|_{p_T^{-1}(t)})$ does not depend on $t\in T$ by the principle of
continuity \cite[Thm. 10.2]{F}. In particular, since $\deg(L_0|_{Y_0})>0$ by Lemma
\ref{deg>0}, we obtain
\begin{equation}\label{deg(L0|Y1)>0}
\deg(L_0|_{Y_1})>0.
\end{equation}



\begin{lemma}\label{ample on PN}
Let $E$ and $pr$ be as in Lemma \ref{deg>0}.

(i) The line bundle $L:=\mathcal{O}_{\mathbb{P}(E^\vee)}(1)\otimes
pr^*\mathcal{O}_{\mathbb{P}^N}(2-\delta_B(E))$ on $\mathbb{P}(E^\vee)$ is ample.

(ii) The line bundle
\begin{equation}\label{sheafAi}
A_i:=L^{r+1}\otimes pr^*\mathcal{O}_{\mathbb{P}^N}(i)\simeq
\mathcal{O}_{\mathbb{P}(E^\vee)}(r+1)\otimes
pr^*\mathcal{O}_{\mathbb{P}^N}((r+1)((2-\delta_B(E))+i)),
\end{equation}
where $r={\rm rk}E$, is also ample for any $i\ge0$.
\end{lemma}
\begin{proof}
(i) We note first that the line bundle $L_0:=L\otimes
pr^*\mathcal{O}_{\mathbb{P}^N}(-1)$ is numerically effective, i.e. the degree of its
restriction to any curve in $\mathbb{P}(E^\vee)$ is positive. Indeed, let $Y$ be an
irreducible curve in $\mathbb{P}(E^\vee)$. If $pr(Y)$ is a curve, then our claim
follows from (\ref{deg(L0|Y1)>0}). If $pr(Y)$ is a point $z$, then $Y\subset
pr^{-1}(z)\simeq\mathbb{P}^{r-1}$ and
$\deg(L_0|_{Y})=\deg(\mathcal{O}_{\mathbb{P}^{r-1}}(1)|_{Y})$ is again positive.

The numerically effective divisor class $c_1(L_0)$ equals $W+(1-\delta_B(E))H$, where
$W:=c_1(\mathcal{O}_{\mathbb{P}(E^\vee)}(1)),\
H:=pr^*c_1(\mathcal{O}_{\mathbb{P}^N}(1))$. By Lemma \ref{ample} the divisor class
$W+a_0H$ on $\mathbb{P}(E^\vee)$ is ample for $a_0-2+\delta_B(E)$ large enough.
Moreover, a corollary of Kleiman's Theorem \cite[Cor. 1.4.9]{L} implies that the
divisor class
$(a_0-2+\delta_B(E))c_1(L_0)+W+a_0H=(a_0-1+\delta_B(E))(W+(2-\delta_B(E))H)$ is ample.
Hence $W+(2-\delta_B(E))H$ is also ample.

(ii) is a direct corollary of (i).
\end{proof}



Recall that a coherent sheaf $\mathcal{F}$ on $\mathbb{P}^N$ is $m-regular$  in the
sense of Mumford-Castelnuovo if $H^i(\mathcal{F}(m-i))=0$ for $i\ge1$.

\begin{theorem}\label{m-regularity}
Let $E$ be a vector bundle of rank $r$ on $\mathbb{P}^N$.

(i) $E$ is $m$-regular for $m\ge m_0:=c_1(E)+(1+r)(2-\delta_B(E))-1$. Furthermore,
$E(m)$ is generated by global sections for $m\ge m_0$.

(ii) For any hyperplane $\mathbb{P}^{N-1}$ in $\mathbb{P}^N$ the vector bundle
$E(m)|_{\mathbb{P}^{N-1}},\ m\ge m_0,$ is generated by global sections and
\begin{equation}\label{h0(E(m0))}
h^0(E(m)|_{\mathbb{P}^{N-1}})\le\frac{r}{(N-1)!} \big(\delta_A(E)+m+N-1\big)^{N-1}.
\end{equation}
\end{theorem}

\begin{proof}
We keep the notations of Lemmas \ref{deg>0} and \ref{ample on PN}. The dualizing sheaf
$\omega_{\mathbb{P}(E^\vee)}$ of $\mathbb{P}(E^\vee)$ is given by the standard formula
\begin{equation}\label{omegaX}
\omega_{\mathbb{P}(E^\vee)}\simeq\mathcal{O}_{\mathbb{P}(E^\vee)}(-r)\otimes
pr^*\mathcal{O}_{\mathbb{P}^N}(c_1(E)-N-1).
\end{equation}
Therefore (\ref{sheafAi}) and (\ref{omegaX}) imply $\omega_{\mathbb{P}(E^\vee)}\otimes
A_i\simeq\mathcal{O}_{\mathbb{P}(E^\vee)}(1)\otimes
pr^*\mathcal{O}_{\mathbb{P}^N}(m_0-N+i)$. Since $A_i$ is ample for $i\ge0$ by Lemma
\ref{ample on PN}, the Kodaira vanishing theorem yields
\begin{equation}\label{Kodaira}
0=H^j(\omega_{\mathbb{P}(E^\vee)}\otimes
A_i)=H^j(\mathcal{O}_{\mathbb{P}(E^\vee)}(1)\otimes
pr^*\mathcal{O}_{\mathbb{P}^N}(m_0-N+i)),\ \ \ i\ge0,\ \ j\ge1.
\end{equation}
In addition, clearly
$$
pr_*(\omega_{\mathbb{P}(E^\vee)}\otimes A_i)\simeq E(m_0-N+i),\ \ \ \
R^jpr_*(\omega_{\mathbb{P}(E^\vee)}\otimes A_i)=0,\ \ j\ge1, i\ge0.
$$
Thus the Leray spectral sequence
$E_2^{aa'}=H^a(R^{a'}pr_*(\omega_{\mathbb{P}(E^\vee)}\otimes A_i)) \Rightarrow
H^{a+a'}(\omega_{\mathbb{P}(E^\vee)}\otimes A_i)$ degenerates and yields (via
(\ref{Kodaira})) $H^j(E(m_0-N+i))=0,\ \ \ i\ge0,\ \ j\ge1$. This shows that $E$ is
$m$-regular for $m\ge m_0$. The fact that, if $E$ is $m$-regular then $E(m)$ is
generated by global sections, is well known \cite[Lem. 1.7.2]{HL}. Assertion (i) is
proved.

Assertion (ii) follows from Le Potier-Simpson's Theorem - see \cite[Lem. 3.3.2]{HL}
and substitute $X=\mathbb{P}^N$, $\deg(X)=1$, $F=E(m)$, $\nu=N-1$,
$X_\nu=\mathbb{P}^{N-1}$, $X_1=\mathbb{P}^1$,
$\mu_{max}(E(m)|_{\mathbb{P}^1})=\delta_A(E(m))=\delta_A(E)+m$.
\end{proof}




\vspace{1cm}

\section{An upper bound for the degree of the variety of maximal jumping lines through a point
of a vector bundle $E$ on $\mathbb{P}^N$}\label{bounded degree}

\vspace{0.5cm}

\begin{sub}{\bf The transformation $L_0$ of a vector bundle $E$ under a linear projection.}
\label{Induction} \rm Let
$$
p_x:\mathbb{P}^N\dashrightarrow\mathbb{P}^{N-1}
$$
be the rational linear projection with center at a point $x\in\mathbb{P}^N$ and let
$\mathbb{\tilde{P}}^N$ be the closure in $\mathbb{P}^N\times\mathbb{P}^{N-1}$ of the
graph of $p_x$. We have the following obvious diagram of projections
\begin{equation}\label{linpr}
\mathbb{P}^N\overset{\sigma}\leftarrow\mathbb{\tilde{P}}^N\overset{\pi}\to\mathbb{P}^{N-1}.
\end{equation}

In this section $E$ will denote a vector bundle of rank $r$ on $\mathbb{P}^N$ with the
additional condition
\begin{equation}\label{dB=0}
\delta_B(E)=0.
\end{equation}
Set $L_0:=\pi_{*}\sigma^*E.$
\begin{theorem}\label{first relations}
(i) $L_0$ is a vector bundle of rank
\begin{equation}\label{rankEN-1}
\rho_0:=\rk L_0=c_1(E)+r
\end{equation}
on $\mathbb{P}^{N-1}$, and its construction is compatible with base change, i.e. for
any cartesian square
\begin{equation}\label{square}
\xymatrix{
\mathcal{X}\ar[d]_{\tilde{\pi}}\ar[r]^{\tilde{g}}&\mathbb{\tilde{P}}^N\ar[d]^{\pi}\\
\mathcal{Y}\ar[r]^g&\mathbb{P}^{N-1}}
\end{equation}
there is a base change isomorphism
$$
\beta_0:\tilde{\pi}_*\tilde{g}^*\sigma^*E \overset{\sim}\to g^*\pi_*\sigma^*E=g^*L_0;
$$
moreover, the natural evaluation map $ev:\pi^*L_0\to\sigma^*E$ is an epimorphism.

(ii) $c_1(L_0)=P(c_1(E),c_2(E))$, where
$P(x,y):=\frac{1}{2}x(x+1)-y\in\mathbb{Q}[x,y]$.

(iii) $\delta_A(E)\ge \delta_A(L_0)$.

(iv) The following inequalities hold:
\begin{equation}\label{ineqty1a}
\delta_A(L_0)\ge-(P(c_1(E),c_2(E)))^2,
\end{equation}
\begin{equation}\label{dBEN-1}
\delta_B(L_0)\ge Q(r,\delta_A(E),c_1(E),c_2(E)),
\end{equation}
where $Q(x,y,z,t):=-(x+z)y+P(z,t)-(P(z,t))^2\in\mathbb{Q}[x,y]$ and the polynomial $P$
is defined in (ii).
\end{theorem}

\begin{proof}
(i) Consider an arbitrary point $y\in\mathbb{P}^{N-1}$ and set
$\mathbb{P}^1_y:=\pi^{-1}(y)$. It follows immediately from (\ref{dB=0}) that
$h^1(E|_{\mathbb{P}^1_y})=0$, hence
$h^0(E|_{\mathbb{P}^1_y})=\chi(E|_{\mathbb{P}^1_y})=c_1(E)+r$. These equalities and
the Base Change Theorem \cite [Ch. 3, Thm. 12.11]{H} imply (\ref{rankEN-1}), the
equality $R^1\pi_*\sigma^*E=0$ and the existence of the isomorphism $\beta_0$.
Moreover, by (\ref{dB=0}) the sheaf $\sigma^*E|_{\mathbb{P}^1_y}$ is generated by
global sections. This means that there is an epimorphism
$ev_y:H^0(\sigma^*E|_{\mathbb{P}^1_y})\otimes\mathcal{O}_{\mathbb{P}^1_y}\twoheadrightarrow
\sigma^*E|_{\mathbb{P}^1_y}$. Moreover, the evaluation map $ev:\pi^*L_0\to\sigma^*E$
is compatible with base change, i.e. we have a commutative diagram
\begin{equation}\label{square with evalmaps}
\xymatrix{ \pi^*L_0\otimes\mathbb{C}_y\ar[d]_{\pi^*\beta_0}
\ar[r]^{\ \ \ \ ev\otimes\mathbb{C}_y}&\sigma^*E\otimes\mathbb{C}_y\ar@{=}[d]\\
H^0(\sigma^*E|_{\mathbb{P}^1_y})\otimes\mathcal{O}_{\mathbb{P}^1_y} \ar@{>>}[r]^{\ \ \
\ \ \ ev_y}&\sigma^*E|_{\mathbb{P}^1_y},}
\end{equation}
where $\pi^*\beta_0$ is an isomorphism; whence the evaluation map
$ev:\pi^*L_0\to\sigma^*E$ is epimorphic.

(ii) For the duration of the proof, fix an arbitrary line
$\mathbb{P}^1\subset\mathbb{P}^{N-1}$ and consider the surface
$S:=\pi^{-1}(\mathbb{P}^1)$ and the projective plane
$\mathbb{P}^2:=\sigma(S)\subset\mathbb{P}^N$. This plane passes through the center $x$
of the projection $p_x:\mathbb{P}^N\dashrightarrow\mathbb{P}^{N-1}$, and diagram
(\ref{linpr}) extends to the commutative diagram
\begin{equation}\label{twosquares}
\xymatrix{ \mathbb{P}^2\ar[d]& S\ar[d]\ar[r]^{\tilde{\pi}}\ar[l]_{\tilde{\sigma}}
&\mathbb{P}^1\ar[d]\\
\mathbb{P}^N&\mathbb{\tilde{P}}^N\ar[r]^{\pi}\ar[l]_{\sigma}&\mathbb{P}^{N-1},}
\end{equation}
where $\tilde{\pi}=\pi|_S$, $\tilde{\sigma}=\sigma|_S:S\to\mathbb{P}^2$ is the blowing
up of $\mathbb{P}^2$ at the point $x$, and the vertical arrows are the inclusions.

Set $\mathcal{O}_S(\tau):=\tilde{\sigma}^*\mathcal{O}_{\mathbb{P}^2}(1),\ \ \
\mathcal{O}_S(h):=\tilde{\pi}^*\mathcal{O}_{\mathbb{P}^1}(1).$ Note that the relations
\begin{equation}\label{tau^2,htau}
\tau^2=\tau h=[pt]
\end{equation}
hold in the Chow ring $A(S)$. Furthermore, $R^1\pi_*\sigma^*E=0$ implies via base
change $R^1\tilde{\pi}_*\tilde{\sigma}^*E|_{\mathbb{P}^2}=0$. Hence Riemann-Roch (see,
e.g., \cite[Ex. 15.2.8]{F}) yields
\begin{equation}\label{Chern1}
{\rm ch}(L_0|_{\mathbb{P}^1})= \tilde{\pi}_*({\rm td}(T_{S/\mathbb{P}^1})\cdot{\rm
ch}(\tilde{\sigma}^*E|_{\mathbb{P}^2})).
\end{equation}
Here $T_{S/\mathbb{P}^1}\simeq\mathcal{O}_S(2\tau-h)$. Therefore, setting
$c_i:=c_i(E),\ i=1,2,$ and using the relations (\ref{tau^2,htau}), we obtain in $A(S)$
\begin{equation}\label{Chern2}
{\rm td}(T_{S/\mathbb{P}^1})\cdot{\rm
ch}(\tilde{\sigma}^*E|_{\mathbb{P}^2})= 1+(r+c_1)\tau-\frac{1}{2}rh+
P(c_1,c_2)[pt].
\end{equation}
In $A(\mathbb{P}^1)$ we have, respectively,
\begin{equation}\label{RR}
\tilde{\pi}_*({\rm td}(T_{S/\mathbb{P}^1})\cdot{\rm
ch}(\tilde{\sigma}^*E|_{\mathbb{P}^2}))= (r+c_1)\cdot1+P(c_1,c_2)[pt].
\end{equation}
Whence (ii) follows.

(iii) Set $\tilde{E}:=E(-\delta_A(E)-1)$. Then $\delta_A(\tilde{E})=-1$ and
\begin{equation}\label{h^0EN=0}
\pi_*\sigma^*\tilde{E}=0,\ \ \ h^0(\tilde{E})=h^0(\sigma^*\tilde{E})=0.
\end{equation}
In addition, (\ref{dB=0}) and the condition $-1=\delta_A(\tilde{E})\ge
\delta_B(\tilde{E})$ give
\begin{equation}\label{d>0}
D:=\delta_A(E)+1=-\delta_B(\tilde{E})>0.
\end{equation}
Let $l_0=\tilde{\sigma}^{-1}(x)$ be the exceptional line on $S$. Since
$\mathcal{O}_S(l_0)=\mathcal{O}_S(\tau-h)$, there is an exact triple
\begin{equation}\label{exact1}
0\to\mathcal{O}_S\to\mathcal{O}_S(D(\tau-h))\to\mathcal{O}_{Dl_0}(Dl_0)\to0.
\end{equation}
Here $Dl_0$ is the standard multiplicity $D$ scheme structure on $l_0$ as a divisor in
$S$, and we have the following exact triples of sheaves of
$\mathcal{O}_{Dl_0}$-modules for $D\ge2$:
$$
0\to\mathcal{O}_{l_0}(-1)\to\mathcal{O}_{Dl_0}(Dl_0)\to\mathcal{O}_{(D-1)l_0}(Dl_0)\to0,
$$
\begin{equation}\label{exact2}
0\to\mathcal{O}_{l_0}(-2)\to\mathcal{O}_{(D-1)l_0}(Dl_0)\to\mathcal{O}_{(D-2)l_0}(Dl_0)\to0,
\end{equation}
$$
.\ .\ .\ .\ .\ .\ .\ .\ .\ .\ .\ .\ .\ .\ .\ .\ .\ .\ .\ .\ .\ .\ .\ .\ .\ .\ .\ .\ .\
.\ .\
$$
$$
0\to\mathcal{O}_{l_0}(1-D)\to\mathcal{O}_{2l_0}(Dl_0)\to\mathcal{O}_{l_0}(-D)\to0.
$$

Note that, similarly to (\ref{h^0EN=0}),
\begin{equation}\label{3)h^0=0}
h^0(\tilde{\sigma}^*\tilde{E}|_{\mathbb{P}^2})=0.
\end{equation}
Next, $l_0=\tilde{\sigma}^{-1}(x)$ implies
$(\tilde{\sigma}^*\tilde{E}|_{\mathbb{P}^2})|_{ml_0}\simeq\mathcal{O}^r_{ml_0}$ for
$m\ge1$. Thus, twisting the triples (\ref{exact2}) by
$\tilde{\sigma}^*\tilde{E}|_{\mathbb{P}^2}$, we obtain for $D\ge2$
\begin{equation}\label{4)h^0=0}
h^0(\tilde{\sigma}^*\tilde{E}|_{\mathbb{P}^2}\otimes\mathcal{O}_{Dl_0}(Dl_0))=0.
\end{equation}
Moreover, (\ref{4)h^0=0}) is evident for $D=1$, hence it holds for $D\ge1$.

Twisting (\ref{exact1}) by $\tilde{\sigma}^*\tilde{E}|_{\mathbb{P}^2}$ we obtain the
exact triple $0\to\tilde{\sigma}^*\tilde{E}|_{\mathbb{P}^2}\to
\tilde{\sigma}^*\tilde{E}|_{\mathbb{P}^2}(D(\tau-h))\to
\tilde{\sigma}^*\tilde{E}|_{\mathbb{P}^2}\otimes\mathcal{O}_{Dl_0}(Dl_0)\to0$, and
(\ref{3)h^0=0}) and (\ref{4)h^0=0}) imply
\begin{equation}\label{5)h^0=0}
0=h^0(\tilde{\sigma}^*\tilde{E}|_{\mathbb{P}^2})(D(\tau-h))=
h^0(\tilde{\pi}_*\tilde{\sigma}^*\tilde{E}|_{\mathbb{P}^2}(D(\tau-h))).
\end{equation}
Applying the base change isomorphism $\beta_0$ to the right square of the diagram
(\ref{twosquares}) and using the projection formula, we get
$\tilde{\pi}_*(\tilde{\sigma}^*\tilde{E}|_{\mathbb{P}^2}(D(\tau-h)))\simeq
(\pi_*\sigma^*\tilde{E}(D)|_{\mathbb{P}^1})(-D)\simeq
(\pi_*\sigma^*E|_{\mathbb{P}^1})(-D)\simeq L_0(-D)|_{\mathbb{P}^1}$. Therefore
(\ref{5)h^0=0}) implies $h^0(L_0(-D)|_{\mathbb{P}^1})=0$, or equivalently
$\delta_A(L_0)< D$ as $\mathbb{P}^1$ is an arbitrary line in $\mathbb{P}^{N-1}$. This
together with (\ref{d>0}) yields (iii).

(iv) Let $L_{0}|_{\mathbb{P}^1}\simeq\underset{i=1}{\overset{\rho_0}\oplus}
\mathcal{O}_{\mathbb{P}^1}(a_i)$. Clearly, $\delta_A(L_0)\ge c_1(L_0)/\rho_0$ as
$a_i\le \delta_A(L_0)$ for $1\le i\le \rho_0$. It is clear also that
$c_1(L_0)/\rho_0\ge-(c_1(L_0))^2$. Therefore $\delta_A(L_0)\ge-(c_1(L_0))^2$. On the
other hand, by (ii), $-(c_1(L_0))^2=-(P(c_1,c_2))^2$. Hence (\ref{ineqty1a}) holds.

Finally, set $\tilde{L}_0:=L_0(-\delta_A(L_0)-1)$. We have
\begin{equation}\label{deltaA= etc}
\delta_A( \tilde{L}_0)=-1,\ \ \ \delta_B(L_0)=\delta_B(\tilde{L}_0)+\delta_A(L_0)+1.
\end{equation}
Assume in addition that the line $\mathbb{P}^1$ in $\mathbb{P}^{N-1}$ is chosen in
such a way that $\delta_B(\tilde{L}_{0}|_{\mathbb{P}^1})=\delta_B(\tilde{L}_0)$. Then
$h^0(\tilde{L}_{0}|_{\mathbb{P}^1})=0$, hence Riemann-Roch yields
\begin{equation}\label{h1(EN-1)}
h^1(\tilde{L}_{0}|_{\mathbb{P}^1})=-c_1(\tilde{L}_0)-\rk\tilde{L}_0=
-c_1(\tilde{L}_0)-\rho_0.
\end{equation}
On the other hand, since $\delta_B(\tilde{L}_0)\le \delta_A(\tilde{L}_0)=-1$, we have
$-\delta_B( \tilde{L}_0)-1=h^1(\mathcal{O}_{\mathbb{P}^1}(\delta_B( \tilde{L}_0))) \le
h^1( \tilde{L}_0|_{\mathbb{P}^1})$. The last two inequalities, together with
(\ref{h1(EN-1)}), imply
\begin{equation}\label{ineqty1b}
-1\ge \delta_B( \tilde{L}_0)\ge c_1( \tilde{L}_0)+\rho_0-1.
\end{equation}
In addition, the definition of $\tilde{L}_0$ and statements (ii) and (iii) imply
$c_1(\tilde{L}_0)+\rho_0-1=-\rho_0(\delta_A(L_0)+1)+c_1(L_0)+\rho_0-1\ge
-\rho_0(\delta_A(E)+1)+P(c_1,c_2)+\rho_0-1.$ Substituting this together with
(\ref{ineqty1b}) into (\ref{deltaA= etc}), and using (\ref{ineqty1a}) and
(\ref{rankEN-1}), we obtain
$$
\delta_B(L_0)\ge-\rho_0\delta_A(E)+P(c_1,c_2)+\delta_A(L_0)
\ge-(r+c_1)\delta_A(E)+P(c_1,c_2)-(P(c_1,c_2))^2,
$$
i.e. (\ref{dBEN-1}).
\end{proof}



\end{sub}

\vspace{0.5cm}

\begin{sub}{\bf An estimate for the transformed kernel of the evaluation map
$\pi^*L_0\to\sigma^*E$.} \label{Ker Eval map} \rm Assume in addition
\begin{equation}\label{integer D}
\delta_A(E)=2\delta,\ \ \ c_1(E)=r\delta
\end{equation}
for some $\delta\in\mathbb{Z}_{>0}$. Set
\begin{equation}\label{gamma}
\gamma:=c_2(E)-\frac{1}{2}r(r-1)\delta^2.
\end{equation}
Then Theorem \ref{first relations} yields
\begin{equation}\label{rho0}
\rho_0=\rk L_0=r(1+\delta),
\end{equation}
\begin{equation}\label{c1L0}
c_1(L_0)=P_1(r,\gamma,\delta),
\end{equation}
\begin{equation}\label{dAL0,dBL0}
Q_1(r,\gamma,\delta)\le \delta_B(L_0)\le \delta_A(L_0)\le2\delta,
\end{equation}
where $P_1(r,\gamma,\delta):=P(r\delta, \gamma+r(r-1)\delta^2/2)$,\ \
$Q_1(r,\gamma,\delta):=Q(r,\ 2\delta,\ r\delta,\ \gamma+r(r-1)\delta^2/2)$.

By Theorem \ref{first relations}(i) we have an exact triple of vector bundles
\begin{equation}\label{exact4}
0\to F\overset{\iota}\to\pi^*L_0\overset{ev}\to\sigma^*E\to0,
\end{equation}
where $F:=\Ker ev$. Restriction to $S$ yields an exact triple
\begin{equation}\label{exact5}
0\to F|_{S}\to\tilde{\pi}^*(L_{0}|_{\mathbb{P}^1})\to\tilde
{\sigma}^*E|_{\mathbb{P}^2}\to0
\end{equation}
and its twisted version
\begin{equation}\label{exact5(kh)}
0\to (F|_{S})(jh)\to\tilde{\pi}^*(L_{0}|_{\mathbb{P}^1}(j))\to\tilde
{\sigma}^*E|_{\mathbb{P}^2}(jh)\to0.
\end{equation}
Base change implies
$L_{0}|_{\mathbb{P}^1}\simeq\tilde{\pi}_*\tilde{\sigma}^*E|_{\mathbb{P}^2}$. Therefore
$L_{0}|_{\mathbb{P}^1}(j)\simeq(\tilde{\pi}_*\tilde{{\sigma}}^*E|_{\mathbb{P}^2})(jh),\
j\in\mathbb{Z}$. Since
$H^1({\tilde{\pi}}^*(L_{0}|_{\mathbb{P}^1}(j)))=H^1(L_{0}|_{\mathbb{P}^1}(j))=0,\ \ \
j\ge-\delta_B(L_0)$, and the morphism $\ H^0(\tilde{\pi}^*(L_{0}|_{\mathbb{P}^1}(j)))
\to H^0(\tilde{{\sigma}}^*E|_{\mathbb{P}^2})(jh)),\ j\in\mathbb{Z}$, induced by
(\ref{exact5(kh)}) is an isomorphism, we obtain
\begin{equation}\label{h1FS(kh)}
h^1(F|_{S}(jh))=0,\ \ \ j\ge-Q_1(r,\gamma,\delta)
\end{equation}
(see (\ref{dAL0,dBL0})).

Next, the triple (\ref{exact5}) implies via (\ref{rho0})-(\ref{c1L0})
$$
\rk F=\rk(F|_{S})=r\delta,
$$
\begin{equation}\label{ciFS}
c_1(F|_{S})=(r\delta)\tau-P_1(r,\gamma,\delta)h,\ \ \
c_2(F|_{S})=P_2(r,\gamma,\delta)[pt],
\end{equation}
where $P_2(r,\gamma,\delta):=-\gamma-r(r-1)\delta^2/2- r^2\delta^2+r\delta
P_1(r,\gamma,\delta)$ and we use the relations (\ref{tau^2,htau}).

Set
\begin{equation}\label{beta}
b:=-\min\{\delta_B(F|_{\mathbb{P}^1_y})\ |\ y\in\mathbb{P}^{N-1}\}
\end{equation}
and observe that $b\ge0$ in view of the monomorphism $\iota$ in (\ref{exact4}). To
obtain an upper bound for $b$, take a point $y\in\mathbb{P}^{N-1}$ such that
$\delta_B(F|_{\mathbb{P}^1_y})=-b$. Then
\begin{equation}\label{F|P1y}
F|_{\mathbb{P}^1_y}\simeq\overset{r\delta}{\underset{i=1}\oplus}
\mathcal{O}_{\mathbb{P}^1}(b_i)=\mathcal{O}_{\mathbb{P}^1}(-b)\oplus
\overset{r\delta}{\underset{i=2}\oplus}\mathcal{O}_{\mathbb{P}^1}(b_i),
\end{equation}
where $-b=b_1\le b_2\le......\le b_{r\delta}\le0$. Restricting (\ref{exact4}) onto
$\mathbb{P}^1_y$ and using (\ref{rho0}), we obtain the triple
$$
0\to F|_{\mathbb{P}^1_y}\to\mathcal{O}_{\mathbb{P}^1_y}^{r(1+\delta)}\to
\sigma^*E|_{\mathbb{P}^1_y}\to0.
$$
Moreover, (\ref{integer D})  yields $\chi(\sigma^*E|_{\mathbb{P}^1_y})=\rk
E+c_1(E)=r(1+\delta)$. Therefore
$0=\chi(F|_{\mathbb{P}^1_y})=-b+\overset{r\delta}{\underset{i=2}\sum}b_i+r(1+\delta)$.
Since $\overset{r\delta}{\underset{i=2}\sum}b_i\le0$, this gives the following upper
bound for $b$:
\begin{equation}\label{beta le}
b=\overset{r\delta}{\underset{i=2}\sum}b_i+ r(1+\delta)\le r(1+\delta).
\end{equation}

Consider the vector bundles
$$
\mathcal{O}_\pi(1):=\sigma^*\mathcal{O}_{\mathbb{P}^N}(1),\ \ \
F_{b}:=F\otimes\mathcal{O}_\pi(b),\ \ \ L_1:=\pi_*F_{b}.
$$
Note that
\begin{equation}\label{R1FN=0}
R^1\pi_*F_{b}=0.
\end{equation}
Furthermore, (\ref{ciFS}) implies
\begin{equation}\label{c1FN|S}
c_1(F_{b}|_S)=r\delta(1+b)\tau-hP_1(r,\gamma,\delta),\ \ \
\end{equation}
\begin{equation}\label{c2FN|S}
c_2(F_{b}|_S)=c_2(F|_S)+(r\delta-1)(r\delta\tau-hP_1(r,\gamma,\delta))b\tau+
\frac{1}{2}rb^2\delta(r\delta-1)[pt].
\end{equation}
Base change, together with (\ref{R1FN=0}), yields
\begin{equation}\label{R1FN|S}
\tilde{\pi}_*(F_{b}|_{S})=L_1|_{\mathbb{P}^1},\ \ \ R^1\tilde{\pi}_*(F_{b}|_{S})=0.
\end{equation}
Hence, by Riemann-Roch (cf. (\ref{Chern1}))
\begin{equation}\label{Chern}
{\rm ch}(L_1|_{\mathbb{P}^1})= \tilde{\pi}_*({\rm td}(T_{S/\mathbb{P}^1})\cdot{\rm
ch}(F_{b}|_{S})).
\end{equation}
Substituting (\ref{c1FN|S}) and (\ref{c2FN|S}) into (\ref{Chern}) and proceeding as in
(\ref{Chern2}) and (\ref{RR}), we obtain
\begin{equation}\label{rkL1}
\rk L_1=r\delta(2+b),
\end{equation}
\begin{equation}\label{c1L1}
c_1(L_1)=F_1(r,b,\gamma,\delta):=\binom{r\delta(1+b)+1}{2}-
(r\delta(1+b)+1)P_1(r,\gamma,\delta)-F(r,b,\gamma,\delta).
\end{equation}
Moreover, (\ref{R1FN|S}) implies
$$
\tilde{\pi}_*(F_{b}|_{S}(jh))=L_1(j)|_{\mathbb{P}^1},\ \ \
R^1\tilde{\pi}_*(F_{b}|_{S}(jh))=0, \ \ \ j\in\mathbb{Z}.
$$
Therefore the Leray spectral sequence
$E^{aa'}_2=H^a(R^{a'}\tilde{\pi}_*(F_{b}|_{S}(jh))\Rightarrow
H^{a+a'}(F_{b}|_{S}(jh))$ degenerates and
\begin{equation}\label{Leray}
H^1(F_{b}|_{S}(jh))=H^1(L_1(j)|_{\mathbb{P}^1}), \ \ \ j\in\mathbb{Z}.
\end{equation}

We are now ready to prove the following lemma.
\begin{lemma}\label{estimate for dABL1}
There exist polynomials $R(x,y,z)$ and $S(x,y,z)$ in $\mathbb{Z}[x,y,z]$ such that
\begin{equation}\label{dAL1,dBL1}
-R(r,\gamma,\delta)\le \delta_B(L_1)\le \delta_A(L_1)\le S(r,\gamma,\delta).
\end{equation}
\end{lemma}
\begin{proof}
Fix a line $l\subset\mathbb{P}^2$ not passing through $x$ (recall that $x$ is the
center of the blow-up $\tilde{\sigma}:S\to\mathbb{P}^2$), and let
$\mathbb{P}^1_{\tau}:=\tilde{\sigma}^{-1}(l)$. Then
$\mathcal{O}_S(\mathbb{P}^1_{\tau})\simeq\mathcal{O}_S(\tau)$. Restricting the triple
(\ref{exact5(kh)}) onto $\mathbb{P}^1_{\tau}$, we obtain an exact triple on
$\mathbb{P}^1_{\tau}\simeq\mathbb{P}^1$
$$
0\to F|_{\mathbb{P}^1_{\tau}}(j)\to
\tilde{\pi}^*(L_0|_{\mathbb{P}^1}(j))|_{\mathbb{P}^1_\tau}\to
\tilde{\sigma}^*(E(j)|_{l})\to0.
$$
Since $\tilde{\pi}^*(L_0|_{\mathbb{P}^1}(j))|_{\mathbb{P}^1_\tau}\simeq
L_0|_{\mathbb{P}^1}(j)$, we have $\chi(E(j)|_{l})=r(1+\delta+j)$. Moreover,
(\ref{rho0}), (\ref{c1L0}) and Riemann-Roch yield $\chi(L_0|_{\mathbb{P}^1}(j))=
P_1(r,\gamma,\delta)+r(1+\delta)(j+1)$. Hence
\begin{equation}\label{chiF(k)|P1tau}
\chi(F|_{\mathbb{P}^1_{\tau}}(j))=P_1(r,\gamma,\delta)+r\delta j.
\end{equation}
Next, (\ref{dAL0,dBL0}) implies
\begin{equation}\label{dAL0+k,dBL0+k}
Q_1(r,\gamma,\delta)+j\le \delta_B(L_0|_{\mathbb{P}^1}(j)) \le
\delta_A(L_0|_{\mathbb{P}^1}(j))\le2\delta+j.
\end{equation}
On the other hand,
$F|_{\mathbb{P}^1_{\tau}}(j)\simeq\overset{r\delta}{\underset{i=1}\oplus}
\mathcal{O}_{\mathbb{P}^1}(e_i),$ where $\delta_B(F|_{\mathbb{P}^1_{\tau}}(j))=e_1\le
e_2\le......\le e_{r\delta}$. Therefore (\ref{chiF(k)|P1tau}) yields
\begin{equation}\label{2chiF(k)|P1tau}
P_1(r,\gamma,\delta)+r\delta j=\chi(F|_{\mathbb{P}^1_{\tau}}(j))=
\delta_B(F|_{\mathbb{P}^1_{\tau}}(j))+r\delta+
\overset{r\delta}{\underset{i=2}\sum}e_i.
\end{equation}

Note that, since $F|_{\mathbb{P}^1_{\tau}}(j)$ is a subbundle of
$\tilde{\pi}^*(L_0|_{\mathbb{P}^1}(j))|_{\mathbb{P}^1_\tau}\simeq
L_0|_{\mathbb{P}^1}(j)$, (\ref{dAL0+k,dBL0+k}) implies $e_2\le......\le e_{r\delta}\le
\delta_A(L_0(j)|_{\mathbb{P}^1})\le2\delta+j$, so that
$\overset{r\delta}{\underset{i=2}\sum}e_i\le(2\delta+j)(r\delta-1)$. This together
with (\ref{2chiF(k)|P1tau}) shows that
$$
\delta_B(F|_{\mathbb{P}^1_{\tau}}(j))=P_1(r,\gamma,\delta)+r\delta(j-1)-
\overset{r\delta}{\underset{i=2}\sum}e_i\ge \delta(r-2r\delta+2\delta)-
P_1(r,\gamma,\delta)+j.
$$
Hence
\begin{equation}\label{dBF|P1tau(k)}
\delta_B(F|_{\mathbb{P}^1_{\tau}}(j))\ge0,\ \ \ j\ge
P_2(r,\gamma,\delta):=P_1(r,\gamma,\delta)+\delta(-r+2r\delta-2\delta),
\end{equation}
and this establishes the implication
\begin{equation}\label{againk ge}
j\ge P_2(r,\gamma,\delta)\ \ \Rightarrow\ \ h^1((F|_{\mathbb{P}^1_{\tau}})(j))=0.
\end{equation}

Consider now the sequence of exact triples
$$
0\to F|_S(jh)\to F|_S(jh+\tau)\to F|_{\mathbb{P}^1_{\tau}}(jh+\tau)\to0,
$$
$$
0\to F|_S(jh+\tau)\to F|_S(jh+2\tau)\to F|_{\mathbb{P}^1_{\tau}}(jh+2\tau)\to0,
$$
$$
\ .\ .\ .\ .\ .\ .\ .\ .\ .\ .\ .\ .\ .\ .\ .\ .\ .\ .\ .\ .\ .\ .\ .\ .\ .\ .\ .\ . \
.\ .\ .\ .\ .\ .\ .\ .
$$
$$
0\to F|_S(jh+(b-1)\tau)\to F|_S(jh+b\tau)\to F|_{\mathbb{P}^1_{\tau}}(jh+b\tau)\to0,
$$
where
\begin{equation}\label{k ge max}
j\ge\max\{-Q_1(r,\gamma,\delta),\ P_2(r,\gamma,\delta)\}.
\end{equation}
Since $\mathcal{O}_S(\tau)|_{\mathbb{P}^1_{\tau}}\simeq
\mathcal{O}_S(\tau)|_{\mathbb{P}^1_{\tau}}\simeq\mathcal{O}_{\mathbb{P}^1_{\tau}}(1)$,
it follows from (\ref{h1FS(kh)}) and (\ref{againk ge}) that
$$
h^1(F|_{S}(jh))=h^1(F|_{\mathbb{P}^1_{\tau}}(jh+i\tau))=0,\ \ 0\le i\le b,
$$
for $j$ as in (\ref{k ge max}). Substituting these equalities subsequently into the
triples in the above sequence and keeping in mind that
$(F|_S)(jh+b\tau)=(F_{b}|_S)(jh)$, we eventually obtain
\begin{equation}\label{h1=0 a}
h^1(F_{b}|_S(jh))=0,\ \ \ j\ge\max\{-Q_1(r,\gamma,\delta), \ P_2(r,\gamma,\delta)\}.
\end{equation}

Set $R(x,y,z):=(-Q_1(x,y,z)^+)+P_2(x,y,z)^+$ (the notation $(\cdot)^+$ is introduced
in section 2). Then (\ref{Leray}) and (\ref{h1=0 a}) imply
$h^1(L_1(j)|_{\mathbb{P}^1})=0,\ \ \ j\ge R(r,\gamma,\delta)$. Hence, since
$\mathbb{P}^1$ is an arbitrary line in $\mathbb{P}^{N-1}$, it follows that
\begin{equation}\label{dBL1}
-R(r,\gamma,\delta)\le \delta_B(L_1).
\end{equation}
This establishes the left-hand side of the inequality (\ref{dAL1,dBL1}).

To obtain the right-hand side, consider a line $\mathbb{P}^1\subset\mathbb{P}^{N-1}$
in diagram (\ref{twosquares}) with $\delta_A(L_1|_{\mathbb{P}^1})=\delta_A(L_1)$ and
$$
L_1|_{\mathbb{P}^1}\simeq\overset{r\delta(2+b)}{\underset{i=1}\oplus}
\mathcal{O}_{\mathbb{P}^1}(a_i),
$$
where $\delta_A(L_1)=a_1\ge a_2\ge...\ge a_{r\delta(2+b)}\ge \delta_B(L_1)$ and $\rk
L_1=r\delta(2+b)$ by (\ref{rkL1}). Note that (\ref{rkL1}), (\ref{c1L1}) and
Riemann-Roch yield
\begin{equation}\label{chi(L1|P1)}
\chi(L_1|_{\mathbb{P}^1})=\rk L_1+c_1(L_1)=r\delta(2+b)+F_1(r,b,\gamma,\delta)
=:F_2(r,b,\gamma,\delta).
\end{equation}
On the other hand, $\chi(L_1|_{\mathbb{P}^1})=\delta_A(L_1)+
\overset{r\delta(2+b)}{\underset{i=2}\sum}a_i$. Whence, in view of (\ref{dBL1}), we
obtain $\delta_A(L_1)=\chi(L_1|_{\mathbb{P}^1})-
\overset{r\delta(2+b)}{\underset{i=2}\sum}a_i
\le\chi(L_1|_{\mathbb{P}^1})-(r\delta(2+b)-1)\delta_B(L_1)
\le\chi(L_1|_{\mathbb{P}^1})+(r\delta(2+b)-1)R(r,\gamma,\delta)$. Combined with
(\ref{chi(L1|P1)}), this yields
\begin{equation}\label{dAL1}
\delta_A(L_1)\le R_1(r,b,\gamma,\delta):= F_2(r,b,\gamma,\delta)+
(r\delta(2+b)-1)R(r,\gamma,\delta).
\end{equation}
Recall that, according to (\ref{beta le}),
\begin{equation}\label{0le beta le}
0\le b\le r(1+\delta).
\end{equation}
Setting $S(x,y,z):=R_1(r,r(1+\delta),\gamma,\delta)^+$, we obtain from (\ref{dAL1})
the desired right-hand side of (\ref{dAL1,dBL1}).

\end{proof}


\end{sub}


\vspace{0.5cm}

\begin{sub}{\bf An estimate for the degree of the variety of maximal jumping lines
$B_\delta^\kappa(E^0,x,\mathbb{P}^{N-1})$.}\label{Deg Jump} \rm Note that
(\ref{F|P1y}) and (\ref{beta}) imply that $F_b|_{\mathbb{P}^1_y}$ is generated by
global sections for any $y\in\mathbb{P}^{N-1}$. Hence base change yields an
epimorphism $\pi^*L_1\twoheadrightarrow F_{b}$ and its twist
\begin{equation}\label{evalmap4}
ev_1:\pi^*L_1\otimes\mathcal{O}_\pi(-b)\twoheadrightarrow F.
\end{equation}
Combining (\ref{exact5}) with (\ref{evalmap4}) we get the exact sequence
\begin{equation}\label{exact6}
\pi^*L_1\otimes\mathcal{O}_\pi(-b)\overset{\Psi}\to\pi^*L_0\to\sigma^*E\to0,
\end{equation}
where $\Psi:=i\circ ev_1$. Twisting (\ref{exact6}) by the $\pi$-relative dualizing
sheaf
$\omega_\pi\simeq\mathcal{O}_\pi(-2)\otimes\pi^*\mathcal{O}_{\mathbb{P}^{N-1}}(1)$,
and applying $R^1\pi_{*}$ we obtain the exact sequence
\begin{equation}\label{exact7}
L_1\otimes A_b\overset{\Phi}\to L_0\to R^1\pi_*(\sigma^*E\otimes\omega_\pi)\to0,
\end{equation}
where
\begin{equation}\label{Ab}
A_b:=(\pi_*\mathcal{O}_\pi(b))^{\vee}\simeq
S^b(\mathcal{O}_{\mathbb{P}^{N-1}}(-1)\oplus\mathcal{O}_{\mathbb{P}^{N-1}})=
\mathcal{O}_{\mathbb{P}^{N-1}}(-b)\oplus\mathcal{O}_{\mathbb{P}^{N-1}}(-b+1)\oplus...
\oplus\mathcal{O}_{\mathbb{P}^{N-1}}.
\end{equation}

Set $E^0:=E(-\delta)$. Then (\ref{integer D}) and (\ref{gamma}) yield
\begin{equation}\label{D,gamma}
c_1(E^0)=0,\ \ \ c_2(E^0)=\gamma,\ \ \ \delta_A(E^0)=\delta=-\delta_B(E^0).
\end{equation}
We set also
\begin{equation}\label{2 Bkd(x)}
B^\kappa_\delta(E^0,x,\mathbb{P}^{N-1}):=\{y\in\mathbb{P}^{N-1}\ |\
\dim_{\mathbb{C}_y}(\mathbb{C}_y\otimes R^1\pi_*(\sigma^*E\otimes\omega_\pi))=\kappa\}
\end{equation}
for $x\in\mathbb{P}^N$.


\vspace{0.4cm} Note that $\kappa\le\rk E=r$. Hence in view of (\ref{rho0}),
\begin{equation}\label{rho0-alpha}
r(1+\delta)=\rho_0\ge \rho_0-\kappa\ge r\delta\ge0.
\end{equation}
Next, denote
\begin{equation}\label{rho1}
\rho_1:=\rk(L_1\otimes A_b)= \rk L_1\cdot\rk A_b=r\delta(2+b)(1+b)
\end{equation}
(we use (\ref{rkL1}) and (\ref{Ab}) here). Observe that (\ref{Ab}) implies
$\delta_A(L_1\otimes A_b)=\delta_A(L_1)-b,\ \ \delta_B(L_1\otimes
A_b)=\delta_B(L_1)-b$, so that \ \
$$
j\delta_B(L_1)-jb=j\delta_B(L_1\otimes A_b)\le \delta_B(\wedge^j(L_1\otimes A_b)) \le
\delta_A(\wedge^j(L_1\otimes A_b)) \le j\delta_A(L_1\otimes A_b)=j\delta_A(L_1)-jb
$$
for any $j\in\mathbb{Z}_{>0}$. This, together with Lemma \ref{estimate for dABL1} and
(\ref{0le beta le}), gives the inequalities
\begin{equation}\label{le dA le}
-jR(r,\gamma,\delta)-jr(1+\delta)\le \delta_B(\wedge^j(L_1\otimes A_b)) \le
\delta_A(\wedge^j(L_1\otimes A_b))\le jS(r,\gamma,\delta).
\end{equation}
In a similar way (\ref{dAL0,dBL0}) gives
\begin{equation}\label{dAwedge p L0}
jQ_1(r,\gamma,\delta)\le j\delta_B(L_0)\le \delta_B(\wedge^jL_0) \le
\delta_A(\wedge^jL_0)\le j\delta_A(L_0)\le2j\delta.
\end{equation}

Notice now that the locally free resolution (\ref{exact7}) of the sheaf
$R^1\pi_*(E\otimes\omega_\pi)$ shows that the $\kappa$-th Fitting ideal sheaf
\footnote{For the definition of Fitting ideals see for instance \cite[p. 492]{E}.} $
\mathcal{I}:=\mathcal{F}itt_{\kappa}(R^1\pi_*(\sigma^*E\otimes\omega_\pi)) $ of the
sheaf $R^1\pi_*(E\otimes\omega_\pi)$ coincides with the image of the morhism
$$
\Lambda:\ \mathcal{E}:= \wedge^{(\rho_0-\kappa)}(L_1\otimes
A_b)\otimes\wedge^{(\rho_0-\kappa)}L_0^{\vee} \to\mathcal{O}_{\mathbb{P}^{N-1}}
$$
induced by the morphism $\Phi$ in (\ref{exact7}). We thus have an epimorphism
\begin{equation}\label{surjn on I}
\mathcal{E}\twoheadrightarrow\mathcal{I}.
\end{equation}
Denote by $V^{\kappa}_\delta(x)$ the subscheme of $\mathbb{P}^{N-1}$ defined by the
ideal sheaf $\mathcal{I}$, i.e.
\begin{equation}\label{3 BKd(x)}
\mathcal{O}_{V^{\kappa}_\delta(x)}:=
\mathcal{O}_{\mathbb{P}^{N-1}}/\mathcal{I}=\coker\Lambda.
\end{equation}
Now (\ref{2 Bkd(x)}) implies
\begin{equation}\label{4 BKd(x)}
B^{\kappa}_\delta(E^0,x,\mathbb{P}^{N-1})={\rm Supp}(\coker\Lambda)
=V^{\kappa}_\delta(x)_{red}.
\end{equation}

Clearly,
$$
\delta_B(\wedge^{(\rho_0-\kappa)}(L_1\otimes A_b))
+\delta_B(\wedge^{(\rho_0-\kappa)}L_0^{\vee})\le \delta_B(\mathcal{E})\le
\delta_A(\mathcal{E})\le \delta_A(\wedge^{(\rho_0-\kappa)}(L_1\otimes A_b))
+\delta_A(\wedge^{(\rho_0-\kappa)}L_0^{\vee}).
$$
Substituting here (\ref{le dA le}) and (\ref{dAwedge p L0}) with $j=\rho_0-\kappa$ and
using (\ref{rho0-alpha}), we obtain
\begin{equation}\label{T12}
T_1(r,\gamma,\delta) \le \delta_B(\mathcal{E})\le \delta_A(\mathcal{E})\le
T_2(r,\gamma,\delta),
\end{equation}
where
$$
T_1(r,\gamma,\delta)=-r(1+\delta)(Q_1(r,\gamma,\delta)-
R(r,\gamma,\delta)-r(1+\delta))^+,
$$
$$
T_2(r,\gamma,\delta)=r(1+\delta)(S(r,\gamma,\delta)+2p\delta)^+.
$$
Furthermore, taking into account (\ref{rho0}) and (\ref{rho1}), we obtain
\begin{equation}\label{rk(E)}
\rk\mathcal{E}=I_0(r,b,\rho_0-\kappa,\delta):=\binom{r\delta(2+b)(1+b)}{\rho_0-\kappa}
\binom{r(1+\delta)}{\rho_0-\kappa}.
\end{equation}
Therefore, using (\ref{0le beta le}) and (\ref{rho0-alpha}) we obtain
\begin{equation}\label{tilde rk(E)}
\rk\mathcal{E}\le I(r,\delta):=I_0(r,r(1+\delta),r(1+\delta),\delta)^+.
\end{equation}
Similarly, (\ref{c1L0}) and (\ref{c1L1}) yield
$$
c_1(\wedge^{(\rho_0-\kappa)}L_0^{\vee})=U_0(r,\gamma,\delta):=P_1(r,\gamma,\delta)
\binom{r(1+\delta)-1}{\rho_0-\kappa-1},
$$
$$
c_1(\wedge^{(\rho_0-\kappa)}L_1\otimes A_b)=U_1(r,b,\gamma,\delta):=
(b+1)\left(F_1(r,b,\gamma,\delta)+ \binom{r(1+\delta)-1}{\rho_0-\kappa-1}r\delta
b(b+1)(b+2)/2\right);
$$
hence
$$
c_1(\mathcal{E})=J_0(r,b,\gamma,\delta):=U_0(r,\gamma,\delta)r\delta(2+b)
(1+b)+U_1(r,\gamma,\delta)r(1+\delta).
$$
Then
\begin{equation}\label{tilde c1(E)}
c_1(\mathcal{E})\le J(r,\gamma,\delta):=J_0(r,r(1+\delta),\gamma,\delta)^+.
\end{equation}

Apply now Theorem \ref{m-regularity} to the bundle $\mathcal{E}$. From (\ref{T12}),
(\ref{tilde rk(E)}) and (\ref{tilde c1(E)}) we obtain that $\mathcal{E}(m_0)$ is
globally generated for
\begin{equation}\label{m0}
m_0=m_0(r,\gamma,\delta):=J(r,\gamma,\delta)+(1+I(r,\delta))(2-T_1(r,\gamma,\delta))-1.
\end{equation}
We thus have an epimorphism
$\mathcal{O}^{t_0}_{\mathbb{P}^{N-1}}\twoheadrightarrow\mathcal{E}(m_0)$, where
\begin{equation}\label{h0}
t_0=t_0(r,\gamma,\delta,N):=I(r,\delta)\big(T_2(r,\gamma,\delta)+m+N-1\big)^{N-1}.
\end{equation}
Hence, by (\ref{surjn on I}), we have an epimorphism
$\mathcal{O}^{t_0}_{\mathbb{P}^{N-1}}\twoheadrightarrow\mathcal{I}(m_0)$. This
epimorphism and the Bezout Theorem show that the degree\footnote{By the degree of a
closed reduced subscheme of $\mathbb{P}^{N-1}$ we mean the sum of degrees of its
irreducible components.} of the reduced closed subscheme
\begin{equation}\label{VKd(x)}
B^{\kappa}_\delta(E^0,x,\mathbb{P}^{N-1})=V^{\kappa}_\delta(x)_{red}.
\end{equation}
of $\mathbb{P}^{N-1}$ satisfies the inequality
$$
\deg B^{\kappa}_\delta(E^0,x,\mathbb{P}^{N-1}) \le\deg V^{\kappa}_\delta(x) \le
m_0^{t_0}.
$$
Substituting here (\ref{m0}) and (\ref{h0}) and using the relations (\ref{D,gamma}),
we obtain the following main result of this section.

\begin{theorem}\label{restricted degree} Let $E^0$ be a rank-$r$
vector bundle on $\mathbb{P}^N$ with $c_1(E^0)=0$,
$\delta_A(E^0)=\delta=-\delta_B(E^0)$ and $\kappa_A(E^0)=\kappa$. Let $l$ be a line in
$\mathbb{P}^N$ with $\delta_A(E^0|_l)=\delta$ and $\kappa_A(E^0|_l)=\kappa$ and let
$x$ be any point on $l$. Let $\mathbb{P}^{N-1}$ be the base of the family of lines
through $x$ in $\mathbb{P}^N$. Then the degree of the reduced closed subscheme
$B^{\kappa}_\delta(E^0,x,\mathbb{P}^{N-1})$ of $\mathbb{P}^{N-1}$ satisfies the
inequality
\begin{equation}\label{degVKd(x)}
\deg B^{\kappa}_\delta(E^0,x,\mathbb{P}^{N-1})\le
m_0(r,c_2(E^0),\delta)^{t_0(r,c_2(E^0),\delta,N)},
\end{equation}
where $m_0(x_1,x_2,x_3)$ and $t_0(x_1,x_2,x_3,x_4)$ are given by (\ref{m0}) and
(\ref{h0}), respectively.
\end{theorem}


\end{sub}


\vspace{1cm}

\section{Proof of Theorem \ref{main}}\label{proof}
\vspace{0.5cm} In the rest of the paper we fix a twisted ind-Grassmannian
$\mathbf{G}=\underset{\to}\lim G(i_m,V^{n_m})$ given by a sequence of embeddings
(\ref{eq1}), and assume that $1<i_m\le n_m-i_m$ for all $m$. We set
$G_m:=G(i_m,V^{n_m})$ and $\tilde{\phi}_m:=\phi_{m-1}\circ...\circ\phi_1$. We fix also
a self-dual vector bundle $\mathbf{E}=\underset{\leftarrow}\lim E_m$ on $\mathbf{G}$
(this means that $E_m\simeq E_m^\vee$ for each $m$) of rank
$\mathbf{r}\in\mathbb{Z}_{>0}$. Then
\begin{equation}\label{d=2dA}
c_1(E_m)=0,\ \ \ \delta(E_m)=2\delta_A(E_m),\ \ \ m\ge1.
\end{equation}

Note that it suffices to prove Theorem \ref{main} for self-dual bundles $\mathbf{E}$.
Indeed, consider an arbitrary finite-rank vector bundle
$\mathbf{E}'=\underset{\leftarrow}\lim E'_m$ on $\mathbf{G}$. Set
$\mathbf{E}=\mathcal{E}nd~\mathbf{E}'$. Since $\mathbf{E}$ is self-dual, we can assume
that Theorem \ref{main} holds for $\mathbf{E}$. Therefore, for any $m$ and any line
$l$ in $G(i_m,V^{n_m})$, $(\mathcal{E}ndE'_m)|_{l}$ is a trivial bundle.
Grothendieck's theorem for vector bundles on $\mathbb{P}^1$ implies immediately that
$E'_m|_l\otimes L_{m,l}$  is a trivial bundle for a suitable line bundle $L_{m,l}$ on $l$.
Since $c_1(L_{m,l})=-c_1(E'_m)/\rk E'_m$ does not depend on $l$, the line bundles $L_{m,l}$
define a line bundle $L_m$ on $G_m$.
Now a standard result in \cite{PT}
(Prop. 1.4.1) shows that $E'_m\otimes L_m$ is trivial for any $m$.
Thus
$\underset{\leftarrow}\lim(E'_m\otimes L_m)$ is trivial.
To see that $\mathbf{E}'$ itself is trivial, note that the line bundles $L_m$ define a line bundle
$\mathbf{L}=\underset{\leftarrow}\lim L_m$ on $\mathbf{G}$. As $\mathbf{G}$ is twisted,
for every $m$ the Chern class $c_1(\mathbf{L}|_{G_m})=c_1(L_m)$ must be divisible by
$\deg(\phi_n\circ...\circ\phi_m)$ for all $n>m$. Since
$\lim_{n\to\infty}\deg(\phi_n\circ...\circ\phi_m)=\infty$, it follows that $c_1(L_m)=0$, and hence $\mathbf{L}$
is trivial. Therefore $\mathbf{E}'\simeq\underset{\leftarrow}\lim(E'_m\otimes L_m)$ is trivial.


\begin{sub}{\bf A first observation on $c_2(\mathbf{E})$.}
\rm Note that the embeddings $\phi_m:G_m\to G_{m+1}$ define homomorphisms
$\phi_m^*:A^2(G_{m+1})\to A^2(G_m)$, and the second Chern class of $\mathbf{E}$ is, by
definition, the projective system $\{c_2(E_m)=\phi_m^*c_2(E_{m+1})\}_{m\ge1}$. Here
$A(G_m)=\underset{i\ge0}\oplus A^i(G_m)$ stands for the Chow ring of $G_m$, and we
recall some standard facts about $A(G_m)$ - cf \cite[14.7]{F}:
\begin{itemize}
\item[(i)] $A^1(G_m)=\Pic(G_m)=\mathbb{Z}[\mathbb{V}_m]$,
$A^2(G_m)=\mathbb{Z}[\mathbb{W}_{1,m}]\oplus\mathbb{Z}[\mathbb{W}_{2,m}]$, where
$\mathbb{\mathbb{V}}_m,\mathbb{W}_{1,m},\mathbb{W}_{2,m}$ are the Schubert varieties
of the form $\mathbb{V}_m=\{V^{i_m}\in G_m|\ \dim(V^{i_m}\cap V_0^{n_m-i_m-1})\ge1$
for a fixed subspace $V_0^{n_m-i_m-1}$ of $V^{n_m}\}$, $\mathbb{W}_{1,m}=\{V^{i_m}\in
G_m| $ $\dim (V^{i_m}\cap V_0^{n_m-i_m-1})\ge1$ for a fixed subspace $V_0^{n_m-i_m-1}$
in $V^{n_m}\}$, $\mathbb{W}_{2,m}=\{{V}^{i_m}\in G_m|\ \dim({V}^{i_m}\cap
V_0^{n_m-i_m+1})\ge2$ for a fixed subspace $V_0^{n_m-i_m+1}$ of $V^{n_m}\}$;
\item[(ii)] $[\mathbb{V}_m]^2=[\mathbb{W}_{1,m}]+[\mathbb{W}_{2,m}]$ in $A^2(G_m)$;
\item[(iii)] there exist integers $a_{ij}(m)\ge0,\ \ i,j=1,2,$ such that
\begin{equation}\label{abcd}
\phi_m^*[\mathbb{W}_{1,m+1}]=a_{11}(m)[\mathbb{W}_{1,m}]+a_{21}(m)[\mathbb{W}_{2,m}],\
\
\phi_m^*[\mathbb{W}_{2,m+1}]=a_{12}(m)[\mathbb{W}_{1,m}]+a_{22}(m)[\mathbb{W}_{2,m}],
\end{equation}
\begin{equation}\label{aij}
\ a_{11}(m)+a_{12}(m)=a_{21}(m)+a_{22}(m)=(\deg\phi_m)^2,\  \ \ m\ge1.
\end{equation}
\end{itemize}

\begin{lemma}\label{limits-of-lambdas}
Given $\mathbf{E}=\underset{\leftarrow}\lim E_m$, there exists an infinite subsequence
of the sequence of Grassmannians $G_m$ such that the coordinates of  $c_2(E_m)$ in the
basis $\{[\mathbb{W}_{1,m}],[\mathbb{W}_{2,m}]\}$ are constants
$\lambda_1\in\mathbb{Z}$ and $\lambda_2\in\mathbb{Z}$. Moreover, if
$\lambda_1\lambda_2\ne0$, then $\lambda_1\lambda_2<0$.
\end{lemma}
\begin{proof}
Let
\begin{equation}\label{c_2(E_m)}
c_2(E_m)=\lambda_{1m}[\mathbb{W}_{1,m}]+\lambda_{2m}[\mathbb{W}_{2,m}].
\end{equation}
Consider the $2\times2$-matrix $A(m)=(a_{ij}(m))$ and the column vector
$\Lambda_m=(\lambda_{1m},\lambda_{2m})^t.$ Relations (\ref{c_2(E_m)}) and (\ref{abcd})
give
\begin{equation}\label{lambdas}
\Lambda_m=A(m)\Lambda_{m+1}.
\end{equation}
Set $\gamma_m:=\lambda_{1m}-\lambda_{2m}$. Then, substituting (\ref{aij}) in
(\ref{lambdas}) we compute
\begin{equation}\label{m to_m+1}
\gamma_m=(a_{11}(m)-a_{21}(m))\gamma_{m+1}=
\gamma_{m+m'+1}\prod_{i=1}^{m'}(a_{11}(m+i)-a_{21}(m+i)),\ \ \ m,m'\ge1.
\end{equation}

Assume that $\gamma_{m_0}\ne0$ for some $m_0\ge1$. Then (\ref{m to_m+1}) implies
$a_{11}(m)-a_{21}(m)\ne0,\ \ \gamma_m\ne0,\ \ m\ge m_0$. Furthermore, if
$|a_{11}(m)-a_{21}(m)|>1$ for an infinite number of values of $m\ge m_0$, then the the
right-hand side of (\ref{m to_m+1}) grows to infinity when $m'\to\infty$, a
contradiction. Hence $|a_{11}(m)-a_{21}(m)|>1$ for at most a finite number of values
of $m\ge m_0$. Removing the Grassmannians $G_m$ with these values of $m$ from our
ind-Grassmannian $\mathbf{G}$ (and taking as new embeddings the corresponding
compositions of old embeddings) we may assume that $|a_{11}(m)-a_{21}(m)|=1$ for all
$m\ge m_0$. Since for an infinite number of values of $m$ the numbers
$a_{11}(m)-a_{21}(m)$ have the same sign, the sequence $\{\gamma_m\}$ has an infinite
constant subsequence. Hence, again by removing appropriate $m$'s in the construction
of $\mathbf{G}$, we may assume
\begin{equation}\label{gamma1}
\gamma:=\gamma_m=\lambda_{1m}-\lambda_{2m}\ne0,\ \ \ m\ge m_0.
\end{equation}

Let $\gamma>0$. (The case $\gamma<0$ is treated similarly.) As it was shown in
\cite[section 5]{PT}, for $m$ large enough, say, for $m\ge m_0$, $\lambda_{1m}$ and
$\lambda_{2m}$ cannot be both nonzero of the same sign. (The argument is carried out
in \cite{PT} for rank-2 bundles but applies to bundles $E_m$ of any rank.) This
property and (\ref{gamma1}) imply that
$$
\gamma\ge\lambda_{1m}\ge0,\ \ 0\ge\lambda_{2m}\ge-\gamma,\ \ m\ge m_0.
$$
Thus, there exist infinite constant subsequences
$\{\lambda_{1,m'}=:\lambda_1\ge0\}_{m'\ge m_0}$ and
$\{\lambda_{2,m'}=:\lambda_2\le0\}_{m'\ge m_0}$, of the sequences
$\{\lambda_{1m}\}_{m\ge m_0}$ and $\{\lambda_{2m}\}_{m\ge m_0}$, respectively. Thus,
again without loss of generality we may assume that the sequences
$\{\lambda_{1m}\}_{m\ge m_0}$ and $\{\lambda_{2m}\}_{m\ge m_0}$ are constant:
\begin{equation}\label{lambda1,2}
0\le\lambda_1=\lambda_{1m},\ \ \ 0\ge\lambda_2=\lambda_{2m},\ \ \ m\ge m_0.
\end{equation}
\end{proof}

In what follows we assume that the coordinates of $c_2(E_m)$ in the basis
$\{[\mathbb{W}_{1,m}],[\mathbb{W}_{2,m}]\}$ are constant for our fixed sequence of
Grassmannians $G_m$.

Recall that there are two families of projective subspaces of maximal dimension in
$G_m$: family I consisting of subspaces $\mathbb{P}^{i_m}=\{V^{i_m}\in G_m\ |\
V^{i_m}\subset V_0^{i_m+1}\},\ \ \ \ \ V_0^{i_m+1}\in G(i_m+1,V^{n_m})$, and family II
consisting of subspaces $\mathbb{P}^{n_m-i_m}=\{V^{i_m}\in G_m\ |\ V^{i_m}\supset
V_0^{i_m-1}\},\ \ \ \ \ V_0^{i_m-1}\in G(i_m-1,V^{n_m})$. Lemma
\ref{limits-of-lambdas} implies therefore the following.

\begin{corollary}\label{cor limits}
In the notations of Lemma \ref{limits-of-lambdas}, we have
$$
c_2(E_m|_{\mathbb{P}^{i_m}})=\lambda_2\ \ \ \ \ {\rm for\ any}\ \mathbb{P}^{i_m}\ {\rm
in\ family}\ {\rm I},
$$
\begin{equation}\label{restrII}
c_2(E_m|_{\mathbb{P}^{n_m-i_m}})=\lambda_1\ \ \ \ \ {\rm for\ any}\
\mathbb{P}^{n_m-i_m}\ {\rm in\ family}\ {\rm II}.
\end{equation}
\end{corollary}

\end{sub}


\vspace{0.5cm}
\begin{sub}{\bf The variety of maximal jumping lines of $E_m$ passing through a point.}
\rm For a fixed $m$, consider the natural diagram
\begin{equation}\label{Gamma(m)}
G_m=G(i_m,V^{n_m})\overset{\pi_1}\leftarrow\Gamma_{m}\overset{\pi_2}\rightarrow
Fl_{m},
\end{equation}
where $\Gamma_{m}:=Fl(i_m-1,i_m,i_m+1,V^{n_m})$, and $Fl_{m}:=Fl(i_m-1,i_m+1,V^{n_m})$
is the base of the family of (projective) lines on $G_m$. Set
$$
Z_a(E_m):=\{l\in Fl_{m}\ |\ \delta_A(E_m|_{l})\ge a\},\ \ \
B_a(E_m):=\pi_2^{-1}(Z_a(E_m)),\ \ \ a\in\mathbb{Z}_{>0}.
$$
The semicontinuity of $\delta_A(E_m|_l)$ as a function of $l$ implies that $Z_a(E_m)$
is closed in $Fl_m$; respectively, $B_a(E_m)$ is closed in $\Gamma_m$. Next, set
\begin{equation}\label{dm}
{\Delta}:=\min\{a\ |\ \im(\pi_1(B_a(E_m))\ne G_m\}-1.
\end{equation}
We then have $Y:=\pi_1(B_{{\Delta}+1}(E_m))\ne G_m,\ \pi_1(B_{\Delta}(E_m))=G_m$, and
$$
G'_m:=G_m\smallsetminus Y= \Bigl\{x\in G_m\ \Big|\ {\Delta}=\max\{\delta_A(E_m|_{l})|\
l\ {\rm is\ a\ line\ on}\ G_m\ {\rm through}\ x\}\Bigr\}
$$
is a dense open subset of $G_m$.

Denote $p_{{\Delta},E_m}:=\pi_1|_{B_{\Delta}(E_m)}$. Then
$B_{\Delta}(E_m)':=p_{{\Delta},E_m}^{-1}(G'_m)$ is closed in $\pi_1^{-1}(G'_m)$ and
the morphism $p_{{\Delta},E_m}:B_{\Delta}(E_m)'\to G'_m$ is projective and surjective.
Similarly, for each $a,\ 1\le a\le \mathbf{r}$,
$$
B_{{\Delta}}^a(E_m):=\{(x,l)\in B_{{\Delta}}(E_m)\ |\ l\in Z_{{\Delta}}(E_m),\ \
\kappa_A(E_m|_l)\ge a\},
$$
is a closed subset in $B_{{\Delta}}(E_m)$; respectively,
$p_{{\Delta},E_m}(B_{{\Delta}}^a(E_m))$ is closed in $G_m$. Since
$\kappa_A(E_m|_l)\ge1$ for any $l\in Z_{{\Delta}}(E_m)$, it follows that
$\pi_1(B_{{\Delta}}^1(E_m))=G_m$.


If $\pi_1(B_{{\Delta}}^{\mathbf{r}}(E_m))\ne G_m$, we put
\begin{equation}\label{km}
K:=\min\{2\le a\le \mathbf{r}\ |\ \pi_1(B_{{\Delta}}^a(E_m))\ne G_m\}-1,\ \ \
T:=\pi_1(B_{{\Delta}}^{K+1}(E_m)),
\end{equation}
\begin{equation}\label{Gm0}
G^0_m:=G'_m\smallsetminus T, \ \ \ B_{{\Delta}}^{K}(E_m)^0:=\pi_1^{-1}(G^0_m)\cap
B_{{\Delta}}^{K}(E_m);
\end{equation}
if $\pi_1(B_{{\Delta}}^{\mathbf{r}}(E_m))=G_m$, we put
\begin{equation}\label{km etc}
K:=\mathbf{r},\ \ \ G^0_m:=G'_m,\ \ \ B_{{\Delta}}^{K}(E_m)^0:=B_{{\Delta}}(E_m).
\end{equation}
By definition, $G^0_m$ is a dense open subset of $G'_m$, hence of $G_m$, and the
morphism
$p_{{\Delta},E_m}^{K}:=\pi_1|_{B_{{\Delta}}^{K}(E_m)^0}:B_{{\Delta}}^{K}(E_m)^0\to
G^0_m$ is projective and surjective.




\vspace{0.5cm}

\end{sub}



\vspace{0.5cm}

\begin{sub}{\bf A bound for the codimension of $B_{\Delta}^K(E_m)$.}\label{codim bounded}
\rm The semicontinuity of $\delta_A(E_m|_l)$ (respectively, of $\delta(E_m|_l)$)
forces the minimal value of of $\delta_A(E_m|_l)$ (respectively, of $\delta(E_m|_l)$)
to be attained on a dense open set of lines $l\in Fl_m$. In what follows we denote
this minimal value by $\delta_A^{gen}(E_m)$ (respectively, by $\delta^{gen}(E_m)$).

\begin{lemma}\label{Descent-Lemma}
Assume $ \delta_A^{gen}(E_m)>\frac{1}{2}\mathbf{r}. $ Then there exists a subsheaf
$\mathcal{F}_m$ of $E_m$ with $c_1(\mathcal{F}_m)>0$.
\end{lemma}
\begin{proof} The inequality $\delta_A^{gen}(E_m)>\frac{1}{2}\mathbf{r}$
and the vanishing of $c_1(E_m)$ imply, for any line $l\subset G_m$ with splitting type
$(\delta_1,...,\delta_\mathbf{r})$ of $E|_l$, that $\delta_s-\delta_{s+1}\ge2$ for
some $s,\ 1\le s\le \mathbf{r}-1$. Thus, the assumption of Theorem 1.4.2 in \cite{PT}
(which is version of the Descent Lemma of \cite[Ch. II, Lemma 2.1.2]{OSS} for a Grassmannian) is
satisfied, and this theorem yields a subsheaf $\mathcal{F}_m$ of $E_m$. Since $E_m$ is
self-dual, the vanishing of $c_1(E_m)$ forces the integer $\delta_s$ to be positive,
hence by the construction of $\mathcal{F}_m$ we have
$c_1(\mathcal{F}_m)=\delta_1+...+\delta_s>0$.
\end{proof}
\begin{lemma}\label{No-subsheaves}
For sufficiently large $m$ there are no subsheaves $\mathcal{F}_m\subset E_m$ with
$c_1(\mathcal{F}_m)>0$.
\end{lemma}
\begin{proof}
Set $\tilde{d}_m:=\deg\tilde{\phi}_m$. Consider the polynomial
$P_m(t):=\tilde{d}_mt+1$ and let
$$
\mathcal{H}_m:=\{C\in\Hilb^{P_m(t)}G_m\ |\ C\ {\rm is\ a\ smooth\ irreducible\
rational\ curve\ of\ degree}\ \tilde{d}_m\ {\rm on}\ G_m\}.
$$
It is well known after Str\o mme [St] that $\mathcal{H}_m$ is a smooth irreducible
variety of dimension $n_m\tilde{d}_m+i_m(n_m-i_m)-3$.

Assume that $\mathcal{F}_m$ is a subsheaf of $E_m$ with $c_1(\mathcal{F}_m)>0$. Then
${\rm codim}_{G_m}\Sing \mathcal{F}_m\ge2$ as $E_m$ is locally free \cite[Ch. II, Cor.
1.1.9]{OSS}. Furthermore, since $G_m$ is a homogeneous space, $\mathcal{H}_m^*:=\{C\in
\mathcal{H}_m\ |\ C\cap{\rm Sing}\mathcal{F}_m=\emptyset\}$ is a dense open subset of
$\mathcal{H}_m$.

Set $a_m:=\underset{C\in\mathcal{H}_m}\min\{\delta_A(E_m|_C\}$. Since
$\delta_A(E_m|_C)$ is semicontinuous as a function of $C$, $\mathcal{H}_m^0:=\{C\in
\mathcal{H}_m\ |\ \delta_A(E_m|_C)=a_m\}$ is a dense open subset of $\mathcal{H}_m$,
and, for any projective line $l\subset G_1$,
\begin{equation}\label{dA(E1)}
\delta_A(E_1|_l)=\delta_A(E_m|_{C_1})\ge a_m,
\end{equation}
where $C_1:=\tilde{\phi}_m(l)\in\mathcal{H}_m$. Now assume that
$c_1(\mathcal{F}_m)\ge1$ and consider any curve
$C\in\mathcal{H}_m^*\cap\mathcal{H}_m^0$ such that $\delta_A(E_m|_C)=a_m$. Since
$\mathcal{F}_m|_C$ is a locally free subsheaf of $E_m|_C$ with
$c_1(\mathcal{F}_m)\ge1$, it follows that $a_m=\delta_A(E_m|_C)\ge
\delta_A(\mathcal{F}_m|_C)\ge
c_1(\mathcal{F}_m|C)/\mathbf{r}=\tilde{d}_mc_1(\mathcal{F}_m)/\mathbf{r}\ge
\tilde{d}_m/\mathbf{r}$. Combining this with (\ref{dA(E1)}) we obtain
$\delta_A(E_1|_l)\ge\tilde{d}_m/\mathbf{r}$, in particular,
$$
b_1:=\underset{l'\in\mathcal{H}_1}\max\{\delta_A(E_1|_{l'})\}\ge\tilde{d}_m/\mathbf{r}.
$$
If $\mathcal{F}_m$ exists for infinitely many values of $m$, the right-hand side of
the last inequality tends to infinity for $m\to\infty$ as $\underset{m\to\infty}\lim
\tilde{d}_m=\infty$, a contradiction.
\end{proof}



\begin{corollary}\label{bounded-jump}
For sufficiently large $m$, \ $\delta_A^{gen}(E_m)\le\frac{1}{2}\mathbf{r}$.
\end{corollary}

Fix $x\in G_m$ and for any $d\in\mathbb{Z}_{>0}$ consider the locally closed subset
$$
B^0_a(x):=\{l\in \pi_1^{-1}(x)\ |\ \delta(E_m|_{l})=a\}
$$
of $\pi_1^{-1}(x)$. Let $B_a(x)$ be its closure in $\pi_1^{-1}(x)$. The semicontinuity
of $\delta(E_m|_{l})$ implies $B^0_a(x)=B_a(x)\smallsetminus(\underset{a'>a}\cup
B_{a'}(x)),\ \ a>0$. Denote
$$
\delta:=\delta_A(E_m),\ \ \ \kappa:=\kappa_A(E_m).
$$
Then
\begin{equation}
B_\delta(x)=B^0_\delta(x).
\end{equation}
Furthermore, put
\begin{equation}\label{Bdk(x)}
B^{\kappa}_\delta(x):=\{l\in B_{\delta}(x)\ |\ \kappa_A(E_m|_l)=\kappa\}
\end{equation}
and note that $B^{\kappa}_\delta(x)$ is a closed subset of $B_{\delta}(x)$. The
following result is proved by A.Tyurin \cite[Ch. 2, \S2, Lemmas 3 and 4]{T}.
\begin{lemma}\label{Tyurin0}
If $B^{\kappa}_\delta(x)\ne\emptyset$, then ${\rm codim}_{\pi_1^{-1}(x)}B\le
\mathbf{r}(\mathbf{r}-1)\delta(E_m)$ for any irreducible component $B$ of
$B^{\kappa}_\delta(x)$.
\end{lemma}
Since $E_m$ is self-dual, it follows that $\delta(E_m)=2\delta$, hence Lemma
\ref{Tyurin0} implies
\begin{equation}\label{2codim}
{\rm codim}_{\pi_1^{-1}(x)}B\le2\mathbf{r}(\mathbf{r}-1)\delta
\end{equation}
whenever $B^{\kappa}_\delta(x)\ne\emptyset$.

Consider the closed subset
$$
W:=\{x\in G_m\ |\ B_{\delta_A^{gen}(E_m)}^0(x)=\emptyset\}
$$
and set
$$
G_m^*:=(G_m\smallsetminus W)\cap G^0_m,
$$
where $G^0_m$ was defined in (\ref{Gm0}) and (\ref{km etc}).

Clearly, $W\cap l=\emptyset$ for a generic line $l\subset G_m$, hence $G_m^*$ is a
dense open subset of $G_m$ and for any $x\in G_m^*$ there exists a line $l\subset G_m$
through $x$ with $\delta_A(E_m|_l)=\delta_A^{gen}(E_m)$.

We need one more result of Tyurin. Lemma 5 in \cite[Ch. 2, \S 2]{T} implies directly
the following.
\begin{corollary}\label{Tyurin2}
There exists a polynomial $F\in\mathbb{Q}[x_1,x_2]$ such that, if $E$ is a self-dual
vector bundle on $\mathbb{P}^3$ and $P$ is an arbitrary plane on $\mathbb{P}^3$, then
$$
\delta_A(E|_l)\le F(\delta^{\rm gen}(E),\chi(E|_{P}))
$$
for any line $l\subset\mathbb{P}^3$.
\end{corollary}

Now fix a point $x\in G_m$ and let $K_m(x)$ be the subvariety of $G_m$ filled by
projective subspaces of maximal dimension in $G_m$ passing through $x$. It is well
known that $K_m(x)$ is a cone over the cartesian product
$\mathbb{P}^{i_m-1}\times\mathbb{P}^{n_m-i_m-1}$. Corollary \ref{Tyurin2} implies
that, for any line $l\in p_1^{-1}(x)$,
\begin{equation}\label{d1|l<poly}
\delta_A(E_m|_{l})\le F(\delta^{\rm gen}(E_m),\chi(E_m|_{P}))
\end{equation}
for some polynomial $F\in\mathbb{Q}[x_1,x_2]$ and some projective plane plane
$P\subset K_m(x)$. The class of $P$ in the Chow ring $A(G_m)$ coincides with the class
of a plane contained in a projective subspace of family I or II. Hence, since
$c_1(E_m)=0$, the Riemann-Roch theorem and Corollary \ref{cor limits} imply that
$\chi(E|_P)$ coincides with $\mathbf{r}-\lambda_1$ or $\mathbf{r}-\lambda_2$.
Substituting this, together with Corollary \ref{bounded-jump}, into (\ref{d1|l<poly})
we see that there exists a constant ${\Delta}$ not depending on $m$ which bounds
$\delta_A(E_{m}|_{l})$ from above for any line $l\in p_1^{-1}(x)$ and any $x\in
G_m^*$.

Passing from the sequence $(G_{m},E_{m})_{m\ge1}$ to its appropriate subsequence
$(G_{m'},E_{m'})_{m'\ge1}$, and replacing the original sequence by this subsequence,
we obtain in view of (\ref{dm}), (\ref{km})-(\ref{km etc}), Lemma \ref{Tyurin0} and
(\ref{2codim}), the following result.
\begin{proposition}\label{subseq}
There exist constants ${\Delta}$, $K$ and $m_0\ge1$ such that for any $m\ge m_0$ there
is a dense open subset $G^*_m$ of $G_m$
such that the following statements hold for any $x\in G^*_m$.

(1) $\delta_A(E_m|_{l})\le {\Delta},\ \ \kappa_A(E_m|_{l})\le K$ for any $l\in
B_m(x)$, and
$$
\delta_A(E_m|_{l})={\Delta},\ \ \ \ \kappa_A(E_m|_{l})=K
$$
for some $l\in B_m(x)$. Therefore, $B^K_{{\Delta}}(x)\ne\emptyset$ and
$$
{\rm codim}_{\pi_1^{-1}(x)}B\le 2\mathbf{r}(\mathbf{r}-1){\Delta}
$$
for any irreducible component $B$ of $B^K_{\Delta}(x)$ according to Lemma
\ref{Tyurin0}.

(2) Set $B_{{\Delta}}^K(E_m)^*:=(p_{{{\Delta}},E_m}^K)^{-1}(G^*_m)$. Then
$p_{{{\Delta}},E_m}^K:B_{{\Delta}}^K(E_m)^*\to G^*_m$ is a projective surjective
morphism such that
\begin{equation}\label{fibre of piDK}
(p_{{{\Delta}},E_m}^K)^{-1}(x)=B_{{\Delta}}^K(x).
\end{equation}
\end{proposition}

\end{sub}




\vspace{0.5cm}

\begin{sub}{\bf Final arguments.}\label{final}
\rm It remains to prove the following.
\begin{theorem}\label{subsheaf-F}
In the framework of Proposition \ref{subseq} assume ${\Delta}>0$. Then there exists a
subsheaf $\mathcal{F}_m$ of $E_m$ with $c_1(\mathcal{F}_m)>0$.
\end{theorem}

\begin{proof}

Consider the relative Grassmannian $g:G(K,E_m)\to G_m$ with fibre
$g^{-1}(x)=G(K,E_m|_x)$ for $x\in G_m$. Set $G(K,E_m)^*:=g^{-1}(G^*_m)$. According to
Proposition \ref{subseq} for any point $(x,l)\in B_{{\Delta}}^K(E_m)^*$ there is a
subbundle
\begin{equation}\label{F(x,l)}
F(x,l)\simeq\mathcal{O}_{\mathbb{P}^1}({\Delta})^K.
\end{equation}
 of $E|_l$. This yields a morphism
\begin{equation}\label{morphism Phi}
\Phi:B_{{\Delta}}^K(E_m)^*\to G(K,E_m)^*,(x,l)\mapsto F(x,l)|_x
\end{equation}
which clearly fits in the commutative diagram
\begin{equation}\label{Phi fits}
\xymatrix{
B_{{\Delta}}^K(E_m)^*\ar[dr]_{p_{{{\Delta}},E_m}^K}\ar[r]^{\Phi} & G(K,E_m)^*\ar[d]^{g}\\
& {\ \ G^*_m}\ .}
\end{equation}

In the rest of the proof we assume that $i_m\ge2$. The remaining case is the case of a
twisted ind-projective space, and we leave it as an exercise to the reader. Note that
(as $i_m\ge2$) $G_m=G(i_m,V^{n_m})$ fits into the diagram (\ref{G(i+1,V)}) for
$V=V^{n_m},\ i=i_m-1$, and set $p:=p_{i_m-1},\ q:=q_{i_m-1}$:
\begin{equation}\label{diagram(10)}
G_m\overset{p}\leftarrow Fl(i_m-1,i_m,V^{n_m})\overset{q}\rightarrow G(i_m-1,V^{n_m}).
\end{equation}
Furthermore, fix a subspace $V^{n_m-1}_0$ in $V^{n_m}$ and put
$Y:=q^{-1}(G(i_m-1,V^{n_m-1}_0))$. The projection $\sigma:=p|_Y:Y\to G_m$ is nothing
but a blow-up of $G_m$ with center at the subvariety
\begin{equation}\label{centerZ}
Z_0:=G(i_m,V^{n_m-1}_0),\ \ \ \ {\rm codim}_{G_m}Z_0=i_m\ge2.
\end{equation}
Fix an arbitrary point $x\in G_m^*\smallsetminus Z_0$ and consider the projective
subspace
\begin{equation}\label{Px}
\mathbb{P}^{n_m-i_m}_x:=\sigma (q^{-1}(q(\sigma^{-1}(x))))\subset G_m
\end{equation}
passing through $x$. Note that the fibre
$B_{{\Delta}}^K(x)=(p_{\Delta,E_m}^K)^{-1}(x)$ of the projection
$p_{\Delta,E_m}^K:B_{{\Delta}}^K(E_m)^*\to G_m^*$ lies in $p_1^{-1}(x)$. Next, setting
$\mathbb{P}^{n_m-i_m-1}(x):=\{\mathbb{P}^{i_m}\ {\rm belongs\ to\ family}\ {\rm I}\ |\
x\in\mathbb{P}^{i_m}\}$ and $\mathbb{P}^{i_m-1}(x):=\{\mathbb{P}^{n_m-i_m}\ {\rm
belongs\ to\ family}\ {\rm II}\ |\ x\in\mathbb{P}^{n_m-i_m}\}$ we obtain the
isomorphism
$$
\mathbb{P}^{i_m-1}(x)\times\mathbb{P}^{n_m-i_m-1}(x)\overset{\simeq}\to
\pi_1^{-1}(x),\ (\mathbb{P}^{n_m-i_m},\mathbb{P}^{i_m})\mapsto
l=\mathbb{P}^{n_m-i_m}\cap\mathbb{P}^{i_m}.
$$
Consider the projections
$$
\mathbb{P}^{n_m-i_m-1}(x)\overset{pr_1}\leftarrow
\pi_1^{-1}(x)\overset{pr_2}\rightarrow \mathbb{P}^{i_m-1}(x).
$$
By the construction of $\sigma$, the base $\mathbb{P}^{n_m-i_m-1}_x$ of the family of
lines through $x$ lying in the subspace $\mathbb{P}^{n_m-i_m}_x$ is a fibre of the
projection $\pi_1^{-1}(x)\overset{pr_2}\rightarrow\mathbb{P}^{i_m-1}(x)$ over a
certain point determined by $x$.

Consider the closed subset
\begin{equation}\label{fibre subset}
B^K_{{\Delta},x}:=B_{{\Delta}}^K(x)\cap\mathbb{P}^{n_m-i_m-1}_x
\end{equation}
in $\mathbb{P}^{n_m-i_m-1}_x$. Proposition \ref{subseq}  implies
\begin{equation}\label{codim bounded 1}
{\rm codim}_{\mathbb{P}^{n_m-i_m-1}_x}X\le2\mathbf{r}(\mathbf{r}-1){{\Delta}}
\end{equation}
for any fixed irreducible component $X$ of $B^K_{{{\Delta}},x}.$ Set
\begin{equation}\label{N=}
N:=2\mathbf{r}(\mathbf{r}-1){{\Delta}}+1.
\end{equation}
Take an arbitrary projective subspace $\mathbb{P}^{N-1}\subset\mathbb{P}^{n_m-i_m-1}_x$ and let
$\mathbb{P}^{N}\subset\mathbb{P}^{n_m-i_m}_x$ be the subspace filled by the lines from
$\mathbb{P}^{N-1}$. Put $E^0:=E_m|_{\mathbb{P}^{N}}$. Then $\delta_A(E^0)={\Delta}$,
$\kappa_A(E^0)=K$ by Proposition \ref{subseq} and $c_2(E^0)=\lambda_1$ by
(\ref{restrII}). In addition, comparing (\ref{2 Bkd(x)}) with (\ref{fibre subset}) and
(\ref{fibre of piDK}), we obtain
$$
B^K_{{\Delta}}(E^0,x,\mathbb{P}^{N-1})=B^K_{{\Delta},x}\cap\mathbb{P}^{N-1},
$$
and (\ref{codim bounded 1}) and (\ref{N=}) imply that
$$
\deg B^K_{{\Delta}}(E^0,x,\mathbb{P}^{N-1})=\deg B^K_{{\Delta},x}
$$
for a generic choice of the subspace $\mathbb{P}^{N-1}$ in $\mathbb{P}^{n_m-i_m-1}_x$.
Applying Theorem \ref{restricted degree} to the vector bundle $E^0$ with
$\delta_A(E^0)={\Delta}$ and $\kappa_A(E^0)=K$, we obtain $\deg
B^{K}_{{\Delta}}(E^0,x,\mathbb{P}^{N-1})\le d$, where $d$ is a constant not depending
on $m$. We thus obtain
\begin{equation}\label{C again}
\deg B^K_{{\Delta},x}\le d.
\end{equation}

Suppose next that $m$ is large enough so that the estimate (\ref{codim bounded 1}) for
any irreducible component of $B^K_{{{\Delta}},x}$ together with the condition
$\underset{m\to\infty}\lim(n_m-i_m)=\infty$ ensure that $B^K_{{{\Delta}},x}$ is
connected. Then $\deg X\le d$ by (\ref{C again}).
We can assume without loss of generality that $\mathbb{P}^{n_m-i_m-1}_x=\Span X$.
Therefore, Theorem \ref{join xx' III} applied to $X$ implies that the following
statement holds. For large enough $m$ any two points of $X$ can be joined by a chain
of subspaces $\mathbb{P}^{\mathbf{k}_0}\subset X$, where $\mathbf{k}_0>\dim
G(K,E_m|_x)$. Thus all such subspaces $\mathbb{P}^{\mathbf{k}_0}$ are mapped by $\Phi$
into the same point. Consequently $\Phi(X)$ is a point, and since $B^K_{{{\Delta}},x}$
is connected, $\Phi(B^K_{{\Delta},x})=\Phi(X)$. This defines a morphism
$$
G^*_m\smallsetminus Z_0\to G(K,E_m|_x), x\mapsto\Phi(B^K_{{\Delta},x}),
$$
hence a subbundle $\mathcal{F}'_m$ of $E_m|_{G^*_m\smallsetminus Z_0}$.


The following well-known construction shows that $\mathcal{F}'_m$ extends to a
subsheaf $\mathcal{F}_m$ of $E_m$. The epimorphism of locally free sheaves
$E_m^\vee|_{G^*_m\smallsetminus Z_0}\twoheadrightarrow (\mathcal{F}'_m)^\vee$ defines
the following composition of embeddings
$\zeta:\mathbb{P}(\mathcal{F}'_m)\hookrightarrow \mathbb{P}(E_m|_{G^*_m\smallsetminus
Z_0})\hookrightarrow \mathbb{P}(E_m)$. Let $U$ be the closure of
$\zeta(\mathbb{P}(\mathcal{F}'_m))$ in $\mathbb{P}(E_m)$. Set
$A:=\mathcal{O}_{\mathbb{P}(E_m)/G_m}(1)$ and let $\theta:\mathbb{P}(E_m)\to G_m$ be
the structure morphism. Applying the functor $R^\cdot\theta_*$ to an exact triple
$0\to\mathcal{I}_{U,\mathbb{P}(E_m)}\otimes A\to A\to A|_U\to0$ we obtain the exact
sequence $E_m^\vee\overset{\epsilon}\to\theta_*(A|_U)\to
R^1\theta_*(\mathcal{I}_{U,\mathbb{P}(E_m)}\otimes A)$. The morphism
$\epsilon|_{G^*_m\smallsetminus Z_0}$ is an epimorphism, hence
$\epsilon^\vee:\mathcal{F}_m:=(\theta_*(A|_U))^\vee\to E_m$ is a monomorphim and
$\mathcal{F}_m|_{G^*_m\smallsetminus Z_0}\simeq \mathcal{F}'_m$.

It remains to show that $c_1(\mathcal{F}_m)>0$. By (\ref{F(x,l)}), if $x\in
G^*_m\smallsetminus Z_0$, for any point $(x,l_0)\in B^K_{{\Delta},x}$ we have a
subbundle $\mathcal{F}(x,l_0)\simeq\mathcal{O}_{\mathbb{P}^1}({\Delta})^K$ of
$E_m|_l$. Hence $c_1(\mathcal{F}(x,l_0))>0$ as ${\Delta}>0$. The line $l_0$ lies in
the projective subspace $\mathbb{P}^{n_m-i_m}_x$ of $G_m$ defined in (\ref{Px}), so
the definition (\ref{fibre subset}) shows that $(y,l_0)\in B^K_{{\Delta},y}$ for any
$y$ in the dense open subset $U:=l_0\cap(G^*_m\smallsetminus Z_0)$ of $l_0$. Therefore
$\mathcal{F}(x,l_0)|_U=\mathcal{F}_m|_U$, and consequently $\mathcal{F}(x,l_0)$ is
isomorphic to a locally free quotient of $\mathcal{F}_m|_{l_0}$, i.e.
$\mathcal{F}(x,l_0)\simeq(\mathcal{F}_m|_{l_0})/{\rm Torsion}(\mathcal{F}_m|_{l_0})$.

Since $\mathcal{F}_m$ is torsion free being a subsheaf of $E_m$, it follows that
$\codim_{G_m}\Sing \mathcal{F}_m\ge2$, so that  $\codim_{Fl_m}\ \pi_2(\pi_1^{-1}(\Sing
\mathcal{F}_m))\ge1$. We thus can find a smooth affine curve $C\subset Fl_m$ with a
marked point $c\in C$ such that $\pi_1(\pi_2^{-1}(c))=l_0$ and
$(C\smallsetminus\{c\})\subset(Fl_m\smallsetminus\pi_2(\pi_1^{-1}(\Sing
\mathcal{F}_m)))$. Consider the ruled surface $S:=\pi_2^{-1}(C)\overset{\pi_2}\to C$
and set $\mathcal{F}_S:=(\pi_1^*\mathcal{F}_m|_S)/{\rm
Torsion}(\pi_1^*\mathcal{F}_m|_S),\ \tilde{l}_0:=\pi_2^{-1}(c),\
\mathcal{F}_{\tilde{l}_0}:=(\mathcal{F}_S|_{\tilde{l}_0})/{\rm
Torsion}(\mathcal{F}_S|_{\tilde{l}_0}),\ \tilde{l}_t:=\pi_2^{-1}(t),\ t\in
C\smallsetminus\{c\}$. The condition $\pi_1(\pi_2^{-1}(c))=l_0$ implies that
$\pi_1|_{\tilde{l}_0}:\tilde{l}_0\to l_0$ is an isomorphism,  hence
$(\pi_1|_{\tilde{l}_0})^*\mathcal{F}(x,l_0)\simeq \mathcal{F}_{\tilde{l}_0}$.
Consequently,
\begin{equation}\label{c_1=A}
c_1(\mathcal{F}(x,l_0))=c_1(\mathcal{F}_{\tilde{l}_0}).
\end{equation}
Furthermore, as $\mathcal{F}_m|_{\pi_1(\tilde{l}_t)}$ is locally free for $t\in
C\smallsetminus\{c\}$ by the inclusion
$(C\smallsetminus\{c\})\subset(Fl_m\smallsetminus\pi_2(\pi_1^{-1}(\Sing
\mathcal{F}_m)))$, it follows that
\begin{equation}\label{c_1=B}
c_1(\mathcal{F}_S|_{\tilde{l}_t})=c_1(\mathcal{F}_m),\ \ \ \ t\in
C\smallsetminus\{c\}.
\end{equation}

We claim that
\begin{equation}\label{c_1 le}
c_1(\mathcal{F}_{\tilde{l}_0})\le c_1(\mathcal{F}_S|_{\tilde{l}_t}),\ \ \ \ t\in
C\smallsetminus\{c\}.
\end{equation}
Indeed, as $\mathcal{F}_S$ is torsion free, using a filtration of $\mathcal{F}_S$ with
rank-1 torsion free consequtive quotients, and removing, if necessary, a finite number
of points from $C\smallsetminus\{c\}$, we reduce the proof of (\ref{c_1 le}) to the
case when $\rk \mathcal{F}_S=1$. Here $\mathcal{F}_S\simeq\mathcal{I}_{Y,S}\otimes L$
for some line bundle $L$ on $S$ and for some subscheme $Y$ of $S$ of dimension $\le0$.
Consider the scheme $Y_0:=Y\cap\tilde{l}_0$ of length $\chi(\mathcal{O}_{Y_0})\ge0$
with support on $\tilde{l}_0$. Then $\mathcal{F}_{\tilde{l}_0}\simeq
L|_{\tilde{l}_0}(-\chi(\mathcal{O}_{Y_0}))$, hence
$c_1(\mathcal{F}_{\tilde{l}_0})=c_1(L|_{\tilde{l}_0})-\chi(\mathcal{O}_{Y_0})=
c_1(L|_{\tilde{l}_t})-\chi(\mathcal{O}_{Y_0})=c_1(\mathcal{F}_S|_{\tilde{l}_t})-\chi(\mathcal{O}_{Y_0})\le
c_1(\mathcal{F}_S|_{\tilde{l}_t}),\ t\in C\smallsetminus\{c\}$.

Finally, (\ref{c_1=A})-(\ref{c_1 le}) imply $c_1(\mathcal{F}_m)\ge
c_1(\mathcal{F}(x,l_0))>0$.
\end{proof}

\begin{corollary}
For all $m>0$ $E_m$ is a trivial vector bundle on $G_m$, and Theorem \ref{main}
follows.
\end{corollary}
\begin{proof} If for sufficiently large $m$, $\Delta>0$, then Theorem \ref{subsheaf-F} contradicts
to Lemma \ref{No-subsheaves}. Hence $\Delta=0$. We are going to show now that this
implies the triviality of $E_m$.

Consider diagram (\ref{diagram(10)}). Note that for any $x\in G_m$ the projective
subspace $\mathbb{P}^{i_m-1}(x)$ of $G(i_m-1,V^{n_m})$ introduced in the proof of
Theorem \ref{subsheaf-F} equals $q(p^{-1}(x))$. Similarly,
$\mathbb{P}^{n_m-i_m}(y):=p(q^{-1}(y))$ is a projective subspace of $G_m$ for $y\in
G(i_m-1,V^{n_m})$. Moreover, it is easy to see that the cone
$K(x):=\underset{y\in\mathbb{P}^{i_m-1}(x)}\cup\mathbb{P}^{n_m-i_m}(y)$ with vertex at
$x$, considered as a reduced subscheme of $G_m$, has the same Zariski tangent space at
$x$ as $G_m$:
\begin{equation}\label{tangent}
T_xK(x)=T_xG_m.
\end{equation}
Since $\Delta=0$, in the notations of the proof of Theorem \ref{subsheaf-F} we have
$B_{{\Delta}}^K(x)=\pi^{-1}(x)$ and $K=\mathbf{r}$ for any $x\in G_m^*$, i.e. $E_m|_l$
is trivial for any projective line $l\subset G_m$ passing through $x$. This implies
that $E_m|_{\mathbb{P}^{n_m-i_m}(y)}$ is trivial for any point
$y\in\mathbb{P}^{i_m-1}(x)$ (see, e.g., \cite[Ch. I, Thm. 3.2.1]{OSS}), and hence
$\mathbb{P}^{n_m-i_m}(y)\subset G_m^*$.

We claim that $G^*_m=G_m$.  Indeed, if $G^*_m\ne G_m$, then for any
$x\in(G_m\smallsetminus G^*_m)$ and any $y\in\mathbb{P}^{i_m-1}(x)$ we have
$\mathbb{P}^{n_m-i_m}(y)\subset(G_m\smallsetminus G^*_m)$. Hence
$K(x)\subset(G_m\smallsetminus G^*_m)$, and by (\ref{tangent}), $T_xG_m=T_xK(x)\subset
T_x(G_m\smallsetminus G^*_m)$, where we consider $(G_m\smallsetminus G^*_m)$ as a
reduced subscheme of $G_m$. Whence $(G_m\smallsetminus G^*_m)=G_m$, contrary to the
fact that $G^*_m$ is a dense open subset of $G_m$.

We have shown that $G^*_m=G_m$ for sufficiently large $m$. As $\Delta=0$, this means
that $E_m|_l$ is trivial for any line $l$ in $G_m$. By \cite[Prop. 1.4.1]{PT} this is
sufficient to conclude that $E_m$ is trivial for large enough $m$, and hence for all
$m>0$.
\end{proof}

\end{sub}




\vspace{2cm}

\begin{thebibliography}{MMMMM}

\bibitem[AK]{AK} A. Altman, S. Kleiman. {\it Foundations of the theory of Fano schemes.}
Compos. math. $\mathbf{34}$ (1977), 3-47.

\bibitem[BM]{BM}R. Braun, S. Muller-Stach. {\it Effective bounds for Nori's connectivity theorem}.
In: Higher dimensional complex varieties, Proc. Int. Conf. Trento, Italy, Walter de
Gruyter 1996, pp. 83�-88.

\bibitem[BV]{BV} W. Barth, A. Van de Ven. {\it On the geometry in codimension 2 in Grassmann
manifolds.} In: Lecture Notes in Math. 412, Springer-Verlag 1974, pp. 1-35.

\bibitem[CT]{CT} I. Coanda, G. Trautmann. {\it The splitting criterion of Kempf and the Babylonian
tower theorem}. Communications of Algebra 34 (2006), 2485-2488.
%
%\bibitem[DiP]{DiP} I. Dimitrov, I. Penkov. {\it Ind-varieties of generalized flags as homogeneous
%spaces for classical ind-groups}. IMRN  No. 55 (2004), 2935-2953.

\bibitem[DP]{DP} J. Donin, I. Penkov. {\it Finite rank vector bundles on inductive limits of
grassmannians}. IMRN No. 34 (2003), 1871-1887.

\bibitem[E]{E} D. Eisenbud. {\it Commutative algebra with a view toward algebraic geometry}.
Springer-Verlag, New York, 1995.

\bibitem[F]{F} W. Fulton. {\it Intersection Theory.} Springer-Verlag, Berlin, 1998.

\bibitem[H]{H} R. Hartshorne. {\it Algebraic Geometry.} Springer-Verlag, New York, 1977.

\bibitem[HL]{HL} D. Huybrechts, M. Lehn. {\it The Geometry of Moduli Spaces of Sheaves}.
Aspects of Math., Vol. E 31, Friedr. Vieweg \&\ Sohn, Braunschweig, 1997.

\bibitem[L]{L} R. Lazarsfeld. {\it Positivity in algebraic geometry,} I�-II.  Ergebnisse der
Mathematik und ihrer Grenzgebiete. 3. Folge, Springer-Verlag, Berlin, 2004.

\bibitem[OSS]{OSS}{ C. Okonek, M. Schneider, H. Spindler.} {\it Vector Bundles on Complex
Projective Spaces.} Birkh\"auser, 1980.

\bibitem[PT]{PT} I. Penkov, A. S. Tikhomirov. {\it Rank 2 vector bundles on ind-Grassmannians}.
Manin's Festschrift, in: Progress in Mathematics 269, Birkhaeuser Boston, 2010, pp. 555-572.

\bibitem[S1]{S1} E. Sato. {\it On the decomposability of infinitely extendable vector bundles on
projective spaces and Grassmann varieties}.  J. Math. Kyoto Univ. $\mathbf{17}$
(1977), 127-150.

\bibitem[S2]{S2} E. Sato. {\it On infinitely extendable vector bundles on $G/P$.} J. Math. Kyoto
Univ. $\mathbf{19}$ (1979), 171-189.

\bibitem[St]{St} S. A. Str\o mme. {\it On parametrized rational curves in Grassmann varieties}.
In: Lecture Notes in Math. 1266,  Springer-Verlag 1987, pp. 251-272.

\bibitem[T]{T} A. N. Tyurin. {\it Vector bundles of finite rank over infinite varieties.} Math.
USSR Izvestija $\mathbf{10}$ (1976), 1187-1204.

\end{thebibliography}




\end{document}
